\documentclass[11pt]{article}

\usepackage[margin=1in]{geometry}
\usepackage{amsmath,amssymb,amsthm}
\usepackage{booktabs}
\usepackage{graphicx}
\usepackage{hyperref}
\usepackage{listings}
\usepackage{array}

\tolerance=9999
\emergencystretch=3em
\hbadness=10000
\raggedbottom

\lstset{
  basicstyle=\ttfamily\footnotesize,
  breaklines=true,
  breakatwhitespace=false,
  columns=fullflexible,
  keepspaces=true,
  mathescape=true,
  frame=single,
  framesep=4pt,
  xleftmargin=2pt,
  xrightmargin=2pt,
  showstringspaces=false,
  tabsize=2,
}

\newtheorem{theorem}{Theorem}
\newtheorem{lemma}{Lemma}
\newtheorem{assumption}{Assumption}
\newtheorem{definition}{Definition}
\newtheorem{remark}{Remark}

\newcommand{\qhat}{\ensuremath{\hat q}}

\title{Supplementary Material\\[2pt]
\large Counting Collapse: A Formally Verified Attractor Theory of\\
Repetition Hallucination in Autoregressive Text-to-Speech}
\author{}
\date{}

\begin{document}
\maketitle

\vspace{-2.5em}
\begin{center}
\emph{This document accompanies the 4-page ICASSP 2026 submission and is
self-contained. It reproduces the formal artifact in full, states precisely
what it does and does not certify, documents the benchmark and measurement
protocol down to the level a re-implementer needs, and reports the negative
result of Section~4.4 of the main paper in the detail it deserves. Every claim
below is traceable to a specific file in the project repository, cited inline
by path.}
\end{center}

\tableofcontents

% ============================================================
\section{The formal development in full}
\label{sec:s1}

The formal artifact lives at \texttt{lean/SpectralTTS/} and is organised into
two files: \texttt{SpectralTTS/CountingCollapse.lean} (Theorem~A and its
supporting lemmas) and \texttt{SpectralTTS/AttentionDilution.lean} (Lemma~B).
Both import all of \texttt{Mathlib} and are checked against
\texttt{leanprover/lean4:v4.34.0-rc1} with \texttt{mathlib} pinned at the same
tag (\texttt{lean/SpectralTTS/lean-toolchain},
\texttt{lean/SpectralTTS/lakefile.toml}).

\paragraph{Typesetting convention.} The listings below reproduce the two Lean
source files exactly: every identifier, every tactic, every line of code and
every word of every doc-comment is preserved verbatim and in its original
order. The one substitution made for typesetting is that Lean's Unicode
mathematical notation is rendered through the corresponding standard \LaTeX{}
math macro, because the base monospace font used in this document does not
contain those glyphs and raw Unicode would silently disappear from the printed
page (we checked this directly before committing to this approach). Concretely,
in the listings that follow: $\forall$ stands for Lean's universal quantifier,
$\exists$ for the existential quantifier, $\le$ for the order relation,
$\to$ for the function-space/implication arrow, $\in$ for set/type
membership, $\sum$ for the finite-sum big operator, $\mathbb{R}$ and
$\mathbb{N}$ for the real and natural number types, $\leftarrow$ for Lean's
``rewrite backwards'' arrow, and $\langle\cdot,\cdot\rangle$ for Lean's
anonymous-constructor brackets; $\delta$ and $\epsilon$ stand for the
one-letter Greek identifiers used throughout both files. Separately,
mathlib's naming convention of suffixing a lemma name with a subscript zero
to mark a positivity/non-degeneracy side condition (rendered here, e.g., as
\texttt{div\_le\_iff0}) is printed with a plain digit rather than a Unicode
subscript, for the same font reason. No token is added, removed, or
reordered.

\subsection{\texttt{CountingCollapse.lean}}
\label{sec:s1-counting}

The file's own header states the scope precisely; we reproduce it first
because it is the most honest single paragraph in the whole artifact:

\begin{lstlisting}
/-
Counting Collapse in Autoregressive Decoders
============================================

Formal core of the paper "Counting Collapse: A Formally Verified Attractor
Theory of Repetition Hallucination in Autoregressive TTS".

Setting.  An autoregressive decoder generating speech for a text prompt that
contains a phrase repeated `k` times passes through a sequence of *repetition
boundary* states `s_0, F s_0, F^2 s_0, ...`, where `F : S $\to$ S` is the composite update
the decoder applies over one full repetition of the phrase.  `S` is the decoder
state space (in practice $\mathbb{R}$^d with the Euclidean metric, which is a nonempty
complete metric space).

The modelling assumption is that `F` is the *same* map at every boundary.  This
is what Lemma B (`AttentionDilution.lean`) buys us: softmax attention over `k`
near-identical repeated text spans cannot tell occurrence `j` from occurrence
`j'`, so the conditioning the decoder receives at each boundary is the same up
to `O($\delta$ + 1/k)`.

What is proved here.  *If* `F` is a contraction (`ContractingWith q F`, i.e.
`q < 1` and `F` is `q`-Lipschitz) then the number of repetitions that any
Lipschitz readout can distinguish is finite and bounded explicitly by
`log(2LC/margin) / log(1/q)`.

What is NOT proved here -- and must not be claimed.  That a real transformer
decoder satisfies the contraction hypothesis.  `q` is an empirical quantity: the
accompanying experiments estimate it as `q$\hat{\ }$` per model from the observed decay of
boundary-state distances.  The Lean artifact certifies the *implication*, not
the premise.
-/
import Mathlib

namespace SpectralTTS

open Filter Function

variable {S : Type*} [MetricSpace S] [CompleteSpace S] [Nonempty S]
variable {F : S $\to$ S} {q : NNReal}
\end{lstlisting}

Every theorem below is stated for a fixed but arbitrary nonempty complete
metric space $S$ (the decoder's hidden-state space, in practice
$\mathbb{R}^d$ with the Euclidean metric) and a fixed map $F : S \to S$ (one
full repetition's worth of decoder update) with contraction modulus $q$. These
are Lean \texttt{variable} declarations, so every subsequent theorem is
implicitly universally quantified over them.

\subsubsection*{\texttt{orbitRadius} and \texttt{dist\_iterate\_fixedPoint\_le}}

\begin{lstlisting}
/-- The a-priori radius of the boundary-state orbit: every iterate of `F`
starting at `s` stays within `orbitRadius F s q * q ^ n` of the fixed point.
Writing it as a named quantity keeps the later bounds readable; it is the `C` of
the paper. -/
noncomputable def orbitRadius (F : S $\to$ S) (s : S) (q : NNReal) : $\mathbb{R}$ :=
  dist s (F s) / (1 - q)

theorem orbitRadius_nonneg (hF : ContractingWith q F) (s : S) :
    0 $\le$ orbitRadius F s q :=
  div_nonneg dist_nonneg hF.one_sub_K_pos.le

/-- **Theorem A(i) -- geometric convergence at repetition boundaries.**
The state after `n` repetitions is within `C q^n` of a single fixed point: the
decoder's record of how many repetitions it has produced decays geometrically.
This is `ContractingWith.apriori_dist_iterate_fixedPoint_le` restated in terms of
`orbitRadius`. -/
theorem dist_iterate_fixedPoint_le (hF : ContractingWith q F) (s : S) (n : $\mathbb{N}$) :
    dist (F^[n] s) (ContractingWith.fixedPoint F hF) $\le$ orbitRadius F s q * (q : $\mathbb{R}$) ^ n := by
  have h := hF.apriori_dist_iterate_fixedPoint_le s n
  calc dist (F^[n] s) (ContractingWith.fixedPoint F hF)
      $\le$ dist s (F s) * (q : $\mathbb{R}$) ^ n / (1 - q) := h
    _ = orbitRadius F s q * (q : $\mathbb{R}$) ^ n := by unfold orbitRadius; ring
\end{lstlisting}

\texttt{orbitRadius} is a bookkeeping device, not new mathematics: it is
literally $C := d(s, Fs)/(1-q)$, the quantity the paper's Theorem~1 calls the
orbit scale. \texttt{orbitRadius\_nonneg} is immediate from $q<1$ (so
$1-q>0$) and $d\ge 0$.

\texttt{dist\_iterate\_fixedPoint\_le} is Theorem~A(i). It is \emph{not}
proved from scratch: mathlib already contains the Banach fixed-point theorem
and, with it, an a-priori error estimate for Picard iteration,
\texttt{ContractingWith.apriori\_dist\_iterate\_fixedPoint\_le}. The classical
argument behind that estimate is the standard telescoping bound: since $F$ is
a $q$-contraction, consecutive iterates satisfy $d(F^i s, F^{i+1}s) \le q^i
d(s,Fs)$, and summing the triangle inequality over $i \ge n$ against the
(unique, mathlib-constructed) fixed point $s^\ast$ gives
$d(F^n s, s^\ast) \le \sum_{i \ge n} q^i\, d(s,Fs) = \frac{q^n}{1-q}\,d(s,Fs)$.
The Lean proof here does no more than \texttt{calc}-rewrite that existing
lemma through the definition of \texttt{orbitRadius}, by \texttt{ring}. The
mathematical content is entirely mathlib's; the paper's contribution at this
step is choosing the right restatement for what follows.

\subsubsection*{\texttt{dist\_iterate\_iterate\_le}}

\begin{lstlisting}
/-- Two repetition counts are separated by at most `C (q^m + q^n)` in state
space.  Both orbits are pinned to the same fixed point, so their mutual distance
inherits the geometric decay. -/
theorem dist_iterate_iterate_le (hF : ContractingWith q F) (s : S) (m n : $\mathbb{N}$) :
    dist (F^[m] s) (F^[n] s) $\le$ orbitRadius F s q * ((q : $\mathbb{R}$) ^ m + (q : $\mathbb{R}$) ^ n) := by
  have hm := dist_iterate_fixedPoint_le hF s m
  have hn := dist_iterate_fixedPoint_le hF s n
  calc dist (F^[m] s) (F^[n] s)
      $\le$ dist (F^[m] s) (ContractingWith.fixedPoint F hF)
        + dist (ContractingWith.fixedPoint F hF) (F^[n] s) := dist_triangle _ _ _
    _ = dist (F^[m] s) (ContractingWith.fixedPoint F hF)
        + dist (F^[n] s) (ContractingWith.fixedPoint F hF) := by rw [dist_comm _ (F^[n] s)]
    _ $\le$ orbitRadius F s q * (q : $\mathbb{R}$) ^ m + orbitRadius F s q * (q : $\mathbb{R}$) ^ n := by gcongr
    _ = orbitRadius F s q * ((q : $\mathbb{R}$) ^ m + (q : $\mathbb{R}$) ^ n) := by ring
\end{lstlisting}

This is the triangle inequality routed through the shared fixed point:
$d(F^m s, F^n s) \le d(F^m s, s^\ast) + d(s^\ast, F^n s) \le Cq^m + Cq^n$, each
half supplied by the previous theorem. \texttt{gcongr} (``generalized
congruence'') is a mathlib tactic that discharges a monotone combination of
already-proved inequalities automatically; here it combines the two
$\le$-facts term by term. This is the state-space statement behind the paper's
claim that any two repetition counts land within $C(q^m+q^n)$ of each other --
it does not yet say anything about a readout.

\subsubsection*{\texttt{exists\_indistinguishable}}

\begin{lstlisting}
/-- **Theorem A(ii) -- eventual indistinguishability.**
For any tolerance `$\epsilon$` there is a horizon `N` beyond which *all* repetition counts
produce states within `$\epsilon$` of one another.  No amount of additional repetition
creates new state-space structure. -/
theorem exists_indistinguishable (hF : ContractingWith q F) (s : S) {$\epsilon$ : $\mathbb{R}$} (h$\epsilon$ : 0 < $\epsilon$) :
    $\exists$ N : $\mathbb{N}$, $\forall$ m n : $\mathbb{N}$, N $\le$ m $\to$ N $\le$ n $\to$ dist (F^[m] s) (F^[n] s) < $\epsilon$ := by
  set C := orbitRadius F s q with hC
  have hC0 : 0 $\le$ C := orbitRadius_nonneg hF s
  have hden : 0 < 2 * C + 1 := by linarith
  obtain $\langle$N, hN$\rangle$ : $\exists$ N : $\mathbb{N}$, (q : $\mathbb{R}$) ^ N < $\epsilon$ / (2 * C + 1) :=
    exists_pow_lt_of_lt_one (by positivity) (by exact_mod_cast hF.1)
  refine $\langle$N, fun m n hm hn => ?_$\rangle$
  have hq0 : (0 : $\mathbb{R}$) $\le$ (q : $\mathbb{R}$) := q.coe_nonneg
  have hq1 : (q : $\mathbb{R}$) $\le$ 1 := le_of_lt (by exact_mod_cast hF.1)
  have hpm : (q : $\mathbb{R}$) ^ m $\le$ (q : $\mathbb{R}$) ^ N := pow_le_pow_of_le_one hq0 hq1 hm
  have hpn : (q : $\mathbb{R}$) ^ n $\le$ (q : $\mathbb{R}$) ^ N := pow_le_pow_of_le_one hq0 hq1 hn
  have hbound := dist_iterate_iterate_le hF s m n
  -- C (q^m + q^n)  $\le$  2C q^N  $\le$  2C $\epsilon$/(2C+1)  <  $\epsilon$
  have h1 : C * ((q : $\mathbb{R}$) ^ m + (q : $\mathbb{R}$) ^ n) $\le$ 2 * C * (q : $\mathbb{R}$) ^ N := by
    have hsum : (q : $\mathbb{R}$) ^ m + (q : $\mathbb{R}$) ^ n $\le$ 2 * (q : $\mathbb{R}$) ^ N := by linarith
    calc C * ((q : $\mathbb{R}$) ^ m + (q : $\mathbb{R}$) ^ n)
        $\le$ C * (2 * (q : $\mathbb{R}$) ^ N) := mul_le_mul_of_nonneg_left hsum hC0
      _ = 2 * C * (q : $\mathbb{R}$) ^ N := by ring
  have h2 : 2 * C * (q : $\mathbb{R}$) ^ N $\le$ 2 * C * ($\epsilon$ / (2 * C + 1)) :=
    mul_le_mul_of_nonneg_left hN.le (by linarith)
  have h3 : 2 * C * ($\epsilon$ / (2 * C + 1)) < $\epsilon$ := by
    have ht : (0 : $\mathbb{R}$) < $\epsilon$ / (2 * C + 1) := div_pos h$\epsilon$ hden
    have heq : $\epsilon$ / (2 * C + 1) * (2 * C + 1) = $\epsilon$ := div_mul_cancel0 $\epsilon$ hden.ne'
    nlinarith [ht, heq]
  linarith [hbound, h1, h2, h3]
\end{lstlisting}

This is Theorem~A(ii). Read operationally, it says: fix any tolerance
$\epsilon>0$; then there is a step count $N$ (a horizon) past which \emph{every
pair} of repetition counts $m,n \ge N$ produces states within $\epsilon$ of
each other. It is the formal counterpart of ``the states the decoder visits
stop being mutually distinguishable,'' the empirical claim tested by the
capacity-gain measurement in the main paper's Section~4.3.

The proof strategy: because $q<1$, $q^N \to 0$, so mathlib's
\texttt{exists\_pow\_lt\_of\_lt\_one} supplies an $N$ with $q^N$ below any
target threshold; we choose the threshold $\epsilon/(2C+1)$ (the ``$+1$'' is
there purely to keep the denominator positive even when $C=0$, i.e.\ when $s$
already \emph{is} the fixed point). For $m,n\ge N$, monotonicity of $q^{(\cdot)}$
on $[0,1]$ gives $q^m,q^n \le q^N$, so by
\texttt{dist\_iterate\_iterate\_le}, $d(F^m s,F^n s) \le C(q^m+q^n) \le 2Cq^N$.
The chain $2Cq^N \le 2C\cdot\frac{\epsilon}{2C+1} < \epsilon$ (the last step
because $\frac{2C}{2C+1}<1$) closes the proof by \texttt{linarith}/\texttt{nlinarith}
on the accumulated facts. Nothing here is specific to attention or to
speech; it is a direct corollary of contraction.

\subsubsection*{\texttt{readout\_gap\_le}}

\begin{lstlisting}
/-- **Theorem A(iii) -- no Lipschitz readout can count.**
Any `L`-Lipschitz map `g : S $\to$ $\mathbb{R}$` -- the model's own stop-token logit, a duration
predictor, or any linear probe -- sees repetition counts `m` and `n` collapsed to
within `L C (q^m + q^n)`. -/
theorem readout_gap_le (hF : ContractingWith q F) {L : NNReal} {g : S $\to$ $\mathbb{R}$}
    (hg : LipschitzWith L g) (s : S) (m n : $\mathbb{N}$) :
    |g (F^[m] s) - g (F^[n] s)| $\le$ (L : $\mathbb{R}$) * (orbitRadius F s q * ((q : $\mathbb{R}$) ^ m + (q : $\mathbb{R}$) ^ n)) := by
  have h1 : |g (F^[m] s) - g (F^[n] s)| = dist (g (F^[m] s)) (g (F^[n] s)) :=
    (Real.dist_eq _ _).symm
  rw [h1]
  calc dist (g (F^[m] s)) (g (F^[n] s))
      $\le$ (L : $\mathbb{R}$) * dist (F^[m] s) (F^[n] s) := hg.dist_le_mul _ _
    _ $\le$ (L : $\mathbb{R}$) * (orbitRadius F s q * ((q : $\mathbb{R}$) ^ m + (q : $\mathbb{R}$) ^ n)) := by
        exact mul_le_mul_of_nonneg_left (dist_iterate_iterate_le hF s m n) L.coe_nonneg
\end{lstlisting}

This is the step that turns a statement about the \emph{state space} into a
statement about \emph{anything a downstream mechanism can compute from the
state}. $g$ is an arbitrary $L$-Lipschitz map from $S$ to $\mathbb{R}$: it
could be the model's own stop-token logit, a linear counting probe trained
after the fact, a duration predictor, or any other bounded-sensitivity
read-out. The proof is two lines: rewrite $|g(a)-g(b)|$ as
$\mathrm{dist}(g(a),g(b))$ on $\mathbb{R}$ (\texttt{Real.dist\_eq}), bound it
by $L\cdot d(a,b)$ (the definition of Lipschitz), then substitute the state-space
bound already proved. No property of $g$ beyond its Lipschitz constant is
used, which is exactly why the theorem applies uniformly to greedy decoding,
sampling, and beam search alike -- the readout argument never inspects how a
token was chosen, only how sensitive a fixed function of the state can be.

\subsubsection*{\texttt{counting\_margin\_le}}

\begin{lstlisting}
/-- **Theorem A(iv) -- the counting horizon is finite (algebraic form).**
If a readout still resolves two repetition counts `m $\le$ n` with decision margin
`margin`, then `q ^ m` cannot have decayed below `margin / (2 L C)`.  Since
`q < 1`, this bounds `m`: only finitely many repetition counts admit a decision
margin at all. -/
theorem counting_margin_le (hF : ContractingWith q F) {L : NNReal} {g : S $\to$ $\mathbb{R}$}
    (hg : LipschitzWith L g) (s : S) {margin : $\mathbb{R}$} {m n : $\mathbb{N}$} (hmn : m $\le$ n)
    (hsep : margin $\le$ |g (F^[m] s) - g (F^[n] s)|) :
    margin $\le$ 2 * (L : $\mathbb{R}$) * orbitRadius F s q * (q : $\mathbb{R}$) ^ m := by
  have hq0 : (0 : $\mathbb{R}$) $\le$ (q : $\mathbb{R}$) := q.coe_nonneg
  have hq1 : (q : $\mathbb{R}$) $\le$ 1 := le_of_lt (by exact_mod_cast hF.1)
  have hpn : (q : $\mathbb{R}$) ^ n $\le$ (q : $\mathbb{R}$) ^ m := pow_le_pow_of_le_one hq0 hq1 hmn
  have hgap := readout_gap_le hF hg s m n
  have hC0 : 0 $\le$ orbitRadius F s q := orbitRadius_nonneg hF s
  have hstep : (L : $\mathbb{R}$) * (orbitRadius F s q * ((q : $\mathbb{R}$) ^ m + (q : $\mathbb{R}$) ^ n))
      $\le$ 2 * (L : $\mathbb{R}$) * orbitRadius F s q * (q : $\mathbb{R}$) ^ m := by
    have hsum : (q : $\mathbb{R}$) ^ m + (q : $\mathbb{R}$) ^ n $\le$ 2 * (q : $\mathbb{R}$) ^ m := by linarith
    have hinner : orbitRadius F s q * ((q : $\mathbb{R}$) ^ m + (q : $\mathbb{R}$) ^ n)
        $\le$ orbitRadius F s q * (2 * (q : $\mathbb{R}$) ^ m) := mul_le_mul_of_nonneg_left hsum hC0
    calc (L : $\mathbb{R}$) * (orbitRadius F s q * ((q : $\mathbb{R}$) ^ m + (q : $\mathbb{R}$) ^ n))
        $\le$ (L : $\mathbb{R}$) * (orbitRadius F s q * (2 * (q : $\mathbb{R}$) ^ m)) :=
          mul_le_mul_of_nonneg_left hinner L.coe_nonneg
      _ = 2 * (L : $\mathbb{R}$) * orbitRadius F s q * (q : $\mathbb{R}$) ^ m := by ring
  linarith
\end{lstlisting}

This is the contrapositive step written out algebraically, before taking
logarithms. Assume the readout \emph{does} separate counts $m\le n$ by at
least \texttt{margin}. Since $m\le n$ and $0\le q\le 1$, $q^n\le q^m$, so
$q^m+q^n \le 2q^m$; feeding this into \texttt{readout\_gap\_le}'s bound
$LC(q^m+q^n)$ gives the tighter bound $2LCq^m$. Chaining
$\texttt{margin} \le LC(q^m+q^n) \le 2LCq^m$ is exactly the statement proved.
In words: whatever separation a Lipschitz readout manages to extract between
counts $m$ and $n$, that separation is itself capped by a quantity that decays
geometrically in the \emph{smaller} of the two counts alone -- which is the
algebraic seed of the horizon bound.

\subsubsection*{\texttt{le\_log\_div\_log\_of\_le\_pow}}

\begin{lstlisting}
/-- Real-arithmetic bridge: a geometric quantity bounded below is bounded in its
exponent.  Used to turn `counting_margin_le` into an explicit horizon. -/
theorem le_log_div_log_of_le_pow {r c : $\mathbb{R}$} (hr0 : 0 < r) (hr1 : r < 1) (hc : 0 < c)
    {m : $\mathbb{N}$} (h : c $\le$ r ^ m) : (m : $\mathbb{R}$) $\le$ Real.log c / Real.log r := by
  have hlogr : Real.log r < 0 := Real.log_neg hr0 hr1
  have hlog : Real.log c $\le$ Real.log (r ^ m) := Real.log_le_log hc h
  rw [Real.log_pow] at hlog
  rw [le_div_iff_of_neg hlogr]
  linarith [hlog]
\end{lstlisting}

A general-purpose real-analysis lemma with no reference to decoders, states,
or contractions: if $0<r<1$, $0<c$, and $c \le r^m$, then
$m \le \log c/\log r$. The subtlety, and the reason this needs its own lemma
rather than a one-line \texttt{gcongr}, is the sign flip: because $0<r<1$,
$\log r < 0$ (\texttt{Real.log\_neg}), so dividing the monotone consequence
$\log c \le m\log r$ (from \texttt{Real.log\_le\_log} applied to $c\le r^m$,
then \texttt{Real.log\_pow} to expand $\log(r^m)=m\log r$) by the
\emph{negative} quantity $\log r$ reverses the inequality's direction --
handled here by mathlib's \texttt{le\_div\_iff\_of\_neg}, which is precisely
the division lemma for a negative denominator. This lemma is the generic
engine that turns \texttt{counting\_margin\_le}'s algebraic inequality into an
explicit bound on $m$ in the next theorem; note that \texttt{scripts/check\_lean.sh}'s
axiom audit does not list this lemma by name (Section~\ref{sec:s1-audit}
explains why that omission is harmless).

\subsubsection*{\texttt{counting\_horizon}}

\begin{lstlisting}
/-- **Theorem A -- Counting Collapse (explicit horizon).**
The headline statement.  Under contraction, a Lipschitz readout with decision
margin `margin > 0` can separate the repetition count `m` from any larger count
only while

    m $\le$ log(margin / (2 L C)) / log q  =  log(2 L C / margin) / log(1 / q).

Beyond that horizon the decoder is provably unable to tell how many repetitions
it has already produced, so its continue/stop decision is decoupled from the
count -- the failure mode observed empirically as looping or truncation. -/
theorem counting_horizon (hF : ContractingWith q F) {L : NNReal} {g : S $\to$ $\mathbb{R}$}
    (hg : LipschitzWith L g) (s : S) {margin : $\mathbb{R}$} {m n : $\mathbb{N}$}
    (hq0 : (0 : $\mathbb{R}$) < (q : $\mathbb{R}$)) (hmargin : 0 < margin) (hmn : m $\le$ n)
    (hsep : margin $\le$ |g (F^[m] s) - g (F^[n] s)|) :
    (m : $\mathbb{R}$) $\le$ Real.log (margin / (2 * (L : $\mathbb{R}$) * orbitRadius F s q)) / Real.log (q : $\mathbb{R}$) := by
  have hq1 : (q : $\mathbb{R}$) < 1 := by exact_mod_cast hF.1
  have hbound := counting_margin_le hF hg s hmn hsep
  have hC0 : 0 $\le$ orbitRadius F s q := orbitRadius_nonneg hF s
  -- a nonzero decision margin forces the prefactor to be strictly positive
  have hprefix : (0 : $\mathbb{R}$) $\le$ 2 * (L : $\mathbb{R}$) * orbitRadius F s q :=
    mul_nonneg (mul_nonneg (by norm_num) L.coe_nonneg) hC0
  have hpos : 0 < 2 * (L : $\mathbb{R}$) * orbitRadius F s q := by
    rcases lt_or_eq_of_le hprefix with h | h
    $\cdot$ exact h
    $\cdot$ exfalso
      rw [$\leftarrow$ h, zero_mul] at hbound
      linarith
  have hratio : margin / (2 * (L : $\mathbb{R}$) * orbitRadius F s q) $\le$ (q : $\mathbb{R}$) ^ m := by
    rw [div_le_iff0 hpos]
    calc margin $\le$ 2 * (L : $\mathbb{R}$) * orbitRadius F s q * (q : $\mathbb{R}$) ^ m := hbound
      _ = (q : $\mathbb{R}$) ^ m * (2 * (L : $\mathbb{R}$) * orbitRadius F s q) := by ring
  exact le_log_div_log_of_le_pow hq0 hq1 (by positivity) hratio
\end{lstlisting}

This is Theorem~A(iv), the operative statement of the paper, combining
\texttt{counting\_margin\_le} with the log bridge just discussed. Two details
are worth pointing out because they are exactly the kind of thing a
pen-and-paper proof would gloss over and Lean forces to be made explicit.

First, before the log bridge can apply, the prefactor $2LC$ must be shown
\emph{strictly} positive (the bridge lemma requires $0<c$ for
$c=\text{margin}/(2LC)$ to be well-formed as a division and for
\texttt{Real.log} to behave). The proof handles the boundary case directly:
if $2LC$ were $0$, then \texttt{hbound} ($\text{margin}\le 2LC\cdot q^m$)
would force $\text{margin}\le 0$, contradicting the hypothesis
$\text{margin}>0$ -- so the very existence of a positive decision margin is
what rules out the degenerate case $L=0$ or $C=0$ (a zero-Lipschitz-constant
readout, or an already-converged orbit). This is the \texttt{rcases
lt\_or\_eq\_of\_le} case split ending in \texttt{exfalso}.

Second, the theorem as stated in Lean gives
$m \le \log\!\big(\text{margin}/(2LC)\big)/\log q$, while the main paper's
Equation~(1) writes the horizon as $\log(2LC/\text{margin})/\log(1/q)$. These
are the same quantity: since $\text{margin}/(2LC) \in (0,1)$ whenever the
readout achieves a real separation (so its log is negative) and $q\in(0,1)$
(so $\log q$ is also negative), $\log(a)/\log(b) = \log(1/a)/\log(1/b)$ for
$a=\text{margin}/(2LC)$, $b=q$ -- a negative-over-negative flips to a
positive-over-positive with the reciprocal arguments. The Lean statement and
the paper's Equation~(1) are the identical horizon, just written with the
sign convention that falls out of \texttt{le\_log\_div\_log\_of\_le\_pow}
versus the convention chosen for exposition.

\subsection{\texttt{AttentionDilution.lean}}
\label{sec:s1-dilution}

\begin{lstlisting}
/-
Attention Dilution
==================

Lemma B of "Counting Collapse: A Formally Verified Attractor Theory of
Repetition Hallucination in Autoregressive TTS".

Setting.  A decoder step attends over the `k` text positions occupied by `k`
repetitions of the same phrase.  Because those positions carry (near-)identical
content, their attention logits are near-identical too: we assume only that the
logits lie within a spread `$\delta$` of one another, which is a directly measurable
quantity.

What is proved.  Every occurrence receives attention weight in
`[e^{-$\delta$}/k, e^{$\delta$}/k]`; consequently the attention profile over the repeated
block is within `(e^{$\delta$} - e^{-$\delta$})/k` of uniform, and its entropy is at least
`log k - $\delta$`.  As `k` grows the profile flattens: the decoder cannot use
attention weight to identify *which* occurrence it is currently rendering.

Why it matters.  This licenses the modelling assumption of
`CountingCollapse.lean` -- that the conditioning at every repetition boundary is
the same map `F`.  The link is quantitative but informal (an `O($\delta$ + 1/k)`
perturbation of the conditioning yields an `O($\delta$ + 1/k)` perturbation of `F`); the
paper states it as a remark, not as a formalized theorem.
-/
import Mathlib

namespace SpectralTTS

open Finset Real

variable {k : $\mathbb{N}$}

/-- Softmax over the `k` attention logits of a repeated text block. -/
noncomputable def softmax (z : Fin k $\to$ $\mathbb{R}$) (i : Fin k) : $\mathbb{R}$ :=
  Real.exp (z i) / $\sum$ j, Real.exp (z j)

theorem sum_exp_pos [NeZero k] (z : Fin k $\to$ $\mathbb{R}$) : 0 < $\sum$ j, Real.exp (z j) :=
  Finset.sum_pos (fun j _ => Real.exp_pos (z j))
    $\langle\langle$0, Nat.pos_of_ne_zero (NeZero.ne k)$\rangle$, mem_univ _$\rangle$

theorem softmax_nonneg [NeZero k] (z : Fin k $\to$ $\mathbb{R}$) (i : Fin k) : 0 $\le$ softmax z i :=
  div_nonneg (Real.exp_pos _).le (sum_exp_pos z).le
\end{lstlisting}

\texttt{softmax} is defined exactly as one would expect --
$\mathrm{softmax}(z)_i = e^{z_i}/\sum_j e^{z_j}$ over the $k$ logits of a
repeated block -- and \texttt{sum\_exp\_pos}/\texttt{softmax\_nonneg} are
routine positivity facts needed to make the subsequent division lemmas
well-formed (the denominator is a finite sum of strictly positive terms over a
nonempty index type, guaranteed by the \texttt{[NeZero k]} instance).

\subsubsection*{\texttt{sum\_softmax}}

\begin{lstlisting}
/-- The softmax profile is a probability distribution. -/
theorem sum_softmax [NeZero k] (z : Fin k $\to$ $\mathbb{R}$) : $\sum$ i, softmax z i = 1 := by
  unfold softmax
  rw [$\leftarrow$ Finset.sum_div, div_self (sum_exp_pos z).ne']
\end{lstlisting}

Establishes that \texttt{softmax} really is a probability distribution over
the $k$ occurrences: $\sum_i e^{z_i}/\sum_j e^{z_j} = \big(\sum_i
e^{z_i}\big)/\sum_j e^{z_j} = 1$, by pulling the common denominator out of the
sum (\texttt{Finset.sum\_div}, used backwards) and cancelling it against
itself. This is a prerequisite for the entropy statement later, since Shannon
entropy is only meaningful over a genuine distribution.

\subsubsection*{\texttt{softmax\_le} and \texttt{le\_softmax}}

\begin{lstlisting}
/-- **Lemma B (upper bound).**  If all logits in the repeated block lie within
`$\delta$` of one another, no single occurrence can capture more than `e^{$\delta$}/k` of the
attention mass. -/
theorem softmax_le [NeZero k] (z : Fin k $\to$ $\mathbb{R}$) {$\delta$ : $\mathbb{R}$} (hz : $\forall$ a b, z a - z b $\le$ $\delta$) (i : Fin k) :
    softmax z i $\le$ Real.exp $\delta$ / k := by
  have hk : (0 : $\mathbb{R}$) < k := Nat.cast_pos.mpr (Nat.pos_of_ne_zero (NeZero.ne k))
  -- every term of the denominator is at least `exp (z i - $\delta$)`
  have hlow : (k : $\mathbb{R}$) * Real.exp (z i - $\delta$) $\le$ $\sum$ j, Real.exp (z j) := by
    have hterm : $\forall$ j $\in$ (univ : Finset (Fin k)), Real.exp (z i - $\delta$) $\le$ Real.exp (z j) := by
      intro j _
      exact Real.exp_le_exp.mpr (by linarith [hz i j])
    calc (k : $\mathbb{R}$) * Real.exp (z i - $\delta$)
        = $\sum$ _j : Fin k, Real.exp (z i - $\delta$) := by
          rw [Finset.sum_const, card_univ, Fintype.card_fin, nsmul_eq_mul]
      _ $\le$ $\sum$ j, Real.exp (z j) := Finset.sum_le_sum hterm
  unfold softmax
  rw [div_le_div_iff0 (sum_exp_pos z) hk]
  have hkey : Real.exp (z i) * (k : $\mathbb{R}$) = Real.exp $\delta$ * ((k : $\mathbb{R}$) * Real.exp (z i - $\delta$)) := by
    rw [Real.exp_sub]
    field_simp
  rw [hkey]
  exact mul_le_mul_of_nonneg_left hlow (Real.exp_pos $\delta$).le

/-- **Lemma B (lower bound).**  Symmetrically, no occurrence can be starved
below `e^{-$\delta$}/k`. -/
theorem le_softmax [NeZero k] (z : Fin k $\to$ $\mathbb{R}$) {$\delta$ : $\mathbb{R}$} (hz : $\forall$ a b, z a - z b $\le$ $\delta$) (i : Fin k) :
    Real.exp (-$\delta$) / k $\le$ softmax z i := by
  have hk : (0 : $\mathbb{R}$) < k := Nat.cast_pos.mpr (Nat.pos_of_ne_zero (NeZero.ne k))
  have hhigh : $\sum$ j, Real.exp (z j) $\le$ (k : $\mathbb{R}$) * Real.exp (z i + $\delta$) := by
    have hterm : $\forall$ j $\in$ (univ : Finset (Fin k)), Real.exp (z j) $\le$ Real.exp (z i + $\delta$) := by
      intro j _
      exact Real.exp_le_exp.mpr (by linarith [hz j i])
    calc $\sum$ j, Real.exp (z j)
        $\le$ $\sum$ _j : Fin k, Real.exp (z i + $\delta$) := Finset.sum_le_sum hterm
      _ = (k : $\mathbb{R}$) * Real.exp (z i + $\delta$) := by
          rw [Finset.sum_const, card_univ, Fintype.card_fin, nsmul_eq_mul]
  unfold softmax
  rw [div_le_div_iff0 hk (sum_exp_pos z)]
  have hkey : Real.exp (-$\delta$) * ((k : $\mathbb{R}$) * Real.exp (z i + $\delta$)) = Real.exp (z i) * (k : $\mathbb{R}$) := by
    rw [Real.exp_add, Real.exp_neg]
    field_simp
  calc Real.exp (-$\delta$) * $\sum$ j, Real.exp (z j)
      $\le$ Real.exp (-$\delta$) * ((k : $\mathbb{R}$) * Real.exp (z i + $\delta$)) :=
        mul_le_mul_of_nonneg_left hhigh (Real.exp_pos _).le
    _ = Real.exp (z i) * (k : $\mathbb{R}$) := hkey
\end{lstlisting}

These two theorems together are Lemma~B's headline bound. The hypothesis
\texttt{hz} says the $k$ logits of the repeated block are mutually within
$\delta$ of each other -- a directly measurable spread, not an assumption
about the mechanism producing the logits.

For the upper bound: since $z_i - z_j \le \delta$ for every $j$, each term
of the denominator satisfies $e^{z_j} \ge e^{z_i-\delta}$, so summing $k$
copies of that lower bound gives $\sum_j e^{z_j} \ge k\, e^{z_i-\delta}$.
Substituting into the definition of softmax,
$\mathrm{softmax}(z)_i = e^{z_i}/\sum_j e^{z_j} \le e^{z_i}/(k\,e^{z_i-\delta})
= e^{\delta}/k$. The lower bound is the mirror argument with the roles of
$i$ and $j$ swapped in the hypothesis ($z_j - z_i \le \delta$, i.e.\
$e^{z_j}\le e^{z_i+\delta}$ for every $j$), giving $\sum_j e^{z_j} \le k\,
e^{z_i+\delta}$ and hence $\mathrm{softmax}(z)_i \ge e^{-\delta}/k$. Both
proofs are the same three-step pattern -- bound every summand, sum $k$
identical bounds via \texttt{Finset.sum\_const}, divide -- executed twice with
the inequality direction and the sign of $\delta$ flipped.

\subsubsection*{\texttt{softmax\_dist\_le}}

\begin{lstlisting}
/-- **Occurrence indistinguishability.**  The attention profile over a block of
`k` repeated spans is within `(e^{$\delta$} - e^{-$\delta$})/k` of uniform.  For fixed `$\delta$`
this vanishes as `1/k`: attention weight carries no information about *which*
repetition is being attended to. -/
theorem softmax_dist_le [NeZero k] (z : Fin k $\to$ $\mathbb{R}$) {$\delta$ : $\mathbb{R}$} (hz : $\forall$ a b, z a - z b $\le$ $\delta$)
    (i j : Fin k) :
    |softmax z i - softmax z j| $\le$ (Real.exp $\delta$ - Real.exp (-$\delta$)) / k := by
  have hsplit : (Real.exp $\delta$ - Real.exp (-$\delta$)) / (k : $\mathbb{R}$)
      = Real.exp $\delta$ / (k : $\mathbb{R}$) - Real.exp (-$\delta$) / (k : $\mathbb{R}$) := by ring
  have hi_up := softmax_le z hz i
  have hj_up := softmax_le z hz j
  have hi_lo := le_softmax z hz i
  have hj_lo := le_softmax z hz j
  rw [abs_sub_le_iff]
  constructor <;> linarith
\end{lstlisting}

A direct corollary: both $\mathrm{softmax}(z)_i$ and $\mathrm{softmax}(z)_j$
lie in the interval $[e^{-\delta}/k,\,e^{\delta}/k]$, an interval of width
$(e^{\delta}-e^{-\delta})/k$, so their difference cannot exceed that width in
either direction (\texttt{abs\_sub\_le\_iff} splits $|x-y|\le c$ into the two
one-sided inequalities $x-y\le c$ and $y-x\le c$, both closed by
\texttt{linarith} from the four bounds already in scope). For fixed $\delta$
this bound is $\Theta(1/k)$: as the block grows, no two occurrences' attention
weights can differ by more than an amount vanishing in $k$ -- the formal
content of ``attention weight carries vanishing information about which
occurrence is being rendered.''

\subsubsection*{\texttt{entropy\_ge}}

\begin{lstlisting}
/-- **Attention entropy grows like `log k`.**  With `p i = softmax z i`, the
Shannon entropy of the profile over the repeated block satisfies
`H(p) $\ge$ log k - $\delta$`: near-uniform attention over `k` identical spans is maximally
uninformative, so the decoder's effective conditioning becomes the block mean and
loses occurrence identity. -/
theorem entropy_ge [NeZero k] (z : Fin k $\to$ $\mathbb{R}$) {$\delta$ : $\mathbb{R}$} (hz : $\forall$ a b, z a - z b $\le$ $\delta$) :
    Real.log k - $\delta$ $\le$ $\sum$ i, softmax z i * (-Real.log (softmax z i)) := by
  have hk : (0 : $\mathbb{R}$) < k := Nat.cast_pos.mpr (Nat.pos_of_ne_zero (NeZero.ne k))
  have key : $\forall$ i : Fin k, (Real.log k - $\delta$) * softmax z i
      $\le$ softmax z i * (-Real.log (softmax z i)) := by
    intro i
    have hup := softmax_le z hz i
    have hlo : 0 < softmax z i := by
      have h1 : Real.exp (-$\delta$) / (k : $\mathbb{R}$) $\le$ softmax z i := le_softmax z hz i
      have h2 : (0 : $\mathbb{R}$) < Real.exp (-$\delta$) / (k : $\mathbb{R}$) := by positivity
      linarith
    have hlog : Real.log (softmax z i) $\le$ $\delta$ - Real.log k := by
      have h1 : Real.log (softmax z i) $\le$ Real.log (Real.exp $\delta$ / k) := Real.log_le_log hlo hup
      rwa [Real.log_div (Real.exp_pos $\delta$).ne' hk.ne', Real.log_exp] at h1
    have hneg : Real.log k - $\delta$ $\le$ -Real.log (softmax z i) := by linarith
    calc (Real.log k - $\delta$) * softmax z i
        $\le$ (-Real.log (softmax z i)) * softmax z i :=
          mul_le_mul_of_nonneg_right hneg hlo.le
      _ = softmax z i * (-Real.log (softmax z i)) := mul_comm _ _
  calc Real.log k - $\delta$
      = (Real.log k - $\delta$) * $\sum$ i, softmax z i := by rw [sum_softmax z]; ring
    _ = $\sum$ i, (Real.log k - $\delta$) * softmax z i := by rw [Finset.mul_sum]
    _ $\le$ $\sum$ i, softmax z i * (-Real.log (softmax z i)) := Finset.sum_le_sum (fun i _ => key i)
\end{lstlisting}

The entropy floor. Write $H(p) = \sum_i p_i(-\log p_i)$ for the Shannon
entropy of the attention profile $p_i = \mathrm{softmax}(z)_i$. The proof
first shows a pointwise inequality (\texttt{key}): for every occurrence $i$,
$(\log k - \delta)\,p_i \le p_i(-\log p_i)$. This follows because
\texttt{softmax\_le} gives $p_i \le e^{\delta}/k$, hence (taking logs, valid
since $p_i>0$ by \texttt{le\_softmax}) $\log p_i \le \delta - \log k$, i.e.\
$-\log p_i \ge \log k - \delta$; multiplying both sides of that by the
nonnegative $p_i$ preserves the inequality. Summing the pointwise bound over
all $k$ occurrences and using $\sum_i p_i = 1$ (\texttt{sum\_softmax}) turns
$\sum_i (\log k-\delta)\,p_i$ into exactly $\log k - \delta$, giving
$H(p) \ge \log k - \delta$. As $k\to\infty$ for fixed $\delta$, the minimum
possible entropy of the attention profile over the repeated block grows
without bound like $\log k$: attention over a large repeated block cannot be
concentrated, no matter how the model's underlying logits are computed,
subject only to their spread being bounded by $\delta$.

\subsection{The axiom audit}
\label{sec:s1-audit}

\texttt{scripts/check\_lean.sh} is the verification gate. It performs three
checks in sequence: (1) \texttt{lake build} succeeds, i.e.\ the project
type-checks against the pinned mathlib revision; (2) \texttt{grep -rn "sorry"}
over \texttt{SpectralTTS/} and \texttt{SpectralTTS.lean} finds no occurrence,
i.e.\ no proof obligation anywhere in the source was left unproved with the
\texttt{sorry} placeholder; and (3) an axiom audit, reproduced verbatim below,
that prints the transitive axiom dependency of eleven exported theorems and
fails the build if the string \texttt{sorryAx} appears anywhere in that
output.

\begin{lstlisting}[mathescape=false]
== axiom audit
cat > /tmp/spectraltts_axioms.lean <<'EOF'
import SpectralTTS
open SpectralTTS
#print axioms SpectralTTS.dist_iterate_fixedPoint_le
#print axioms SpectralTTS.dist_iterate_iterate_le
#print axioms SpectralTTS.exists_indistinguishable
#print axioms SpectralTTS.readout_gap_le
#print axioms SpectralTTS.counting_margin_le
#print axioms SpectralTTS.counting_horizon
#print axioms SpectralTTS.softmax_le
#print axioms SpectralTTS.le_softmax
#print axioms SpectralTTS.softmax_dist_le
#print axioms SpectralTTS.entropy_ge
#print axioms SpectralTTS.sum_softmax
EOF
OUT=$(lake env lean /tmp/spectraltts_axioms.lean 2>&1)
echo "$OUT"
if echo "$OUT" | grep -q "sorryAx"; then
  echo "FAIL: sorryAx in dependency closure"; exit 1
fi
\end{lstlisting}

\paragraph{What the three admitted axioms mean.} Lean's kernel tracks, for
every proof term, exactly which axioms it ultimately rests on beyond pure
type-theoretic reduction. Three axioms are considered standard and
non-suspicious throughout ordinary (i.e.\ classical) formalised mathematics,
and mathlib is built expecting all three:
\begin{itemize}
\item \textbf{\texttt{propext}} (propositional extensionality): two
  propositions that are provably logically equivalent are equal as terms of
  type \texttt{Prop}. This licenses rewriting one true statement for another
  and is used pervasively any time an equality between propositions, rather
  than merely an if-and-only-if, is needed internally.
\item \textbf{\texttt{Classical.choice}}: the axiom of choice, in the form
  that every nonempty type has a term. Mathlib is a classical (not
  constructive) library and uses this axiom routinely -- for instance, to
  decide the truth of an arbitrary proposition about real numbers
  (\texttt{Real.log}'s case splits, order comparisons, and the existence
  statement in \texttt{ContractingWith.fixedPoint} all rest on classical
  reasoning of this kind).
\item \textbf{\texttt{Quot.sound}}: quotient soundness. Lean 4 builds
  quotient types (used, transitively, to construct $\mathbb{R}$ as a quotient
  of Cauchy sequences, \texttt{NNReal} as a subtype of that, and
  \texttt{Finset} as a quotient of \texttt{Multiset} which is itself a
  quotient of \texttt{List}) and this axiom is what makes the definitional
  behaviour of quotients sound at the kernel level.
\end{itemize}
Essentially every nontrivial mathlib theorem depends on some subset of these
three; a formal result is treated as unconditionally established in this
ecosystem exactly when its axiom dependency is a subset of
$\{\texttt{propext}, \texttt{Classical.choice}, \texttt{Quot.sound}\}$.

\paragraph{What \texttt{sorryAx} would mean.} Whenever a proof anywhere in a
term's transitive dependency closure contains an unproved hole -- written
\texttt{sorry}, or produced by an incomplete tactic block -- Lean's elaborator
inserts a synthetic axiom called \texttt{sorryAx} standing in for ``trust me.''
A theorem whose dependency closure contains \texttt{sorryAx} is not actually
proved: it is, formally, an assumption. Because \texttt{sorryAx} propagates
through every downstream use of the offending lemma, a single unproved hole
anywhere upstream -- even inside a helper three calls deep, even hypothetically
inside a buggy mathlib lemma -- would surface in the audit of every theorem
that transitively depends on it. This is why the \texttt{grep -rn "sorry"}
textual scan (step 2) and the semantic \texttt{\#print axioms} scan (step 3)
are both run: the grep catches \texttt{sorry} written directly in this
project's own two files, while the axiom audit is the only check that would
also catch a hole hidden inside a project-local tactic-generated term that the
textual scan cannot see, or (in principle) an inconsistency imported from
mathlib itself. Passing both certifies that all eleven audited theorems
reduce, through Lean's kernel type-checker, to a proof term whose only axioms
are the three standard ones above.

\paragraph{Coverage note.} The audit list contains eleven theorems: the seven
from \texttt{CountingCollapse.lean} covered in
Section~\ref{sec:s1-counting} except \texttt{le\_log\_div\_log\_of\_le\_pow},
plus the four listed and named theorems from
\texttt{AttentionDilution.lean} (\texttt{softmax\_le}, \texttt{le\_softmax},
\texttt{softmax\_dist\_le}, \texttt{entropy\_ge}) and \texttt{sum\_softmax}.
\texttt{le\_log\_div\_log\_of\_le\_pow} is not itself named in the
\texttt{\#print axioms} list, but this is not a coverage gap: it is a direct
dependency of \texttt{counting\_horizon}, so Lean's kernel necessarily
re-checks it (and everything \emph{it} depends on) while type-checking
\texttt{counting\_horizon}, and any axiom it introduced would appear in
\texttt{counting\_horizon}'s own reported axiom set. Auditing
\texttt{counting\_horizon} therefore already certifies
\texttt{le\_log\_div\_log\_of\_le\_pow} transitively. The same reasoning
covers the unlisted helper lemmas \texttt{orbitRadius\_nonneg},
\texttt{sum\_exp\_pos}, and \texttt{softmax\_nonneg}: each is a dependency of
an audited theorem.

\paragraph{What is and is not proved (restated plainly).} The Lean artifact
certifies an implication, not a fact about trained transformers.
Concretely: \emph{if} a decoder's per-repetition update map $F$ is a genuine
$q$-contraction ($q<1$, Lipschitz) on its hidden-state space, \emph{then}
Theorem~A(i)--(iv) follow with no further assumptions, machine-checked down to
three standard axioms. Whether any given trained decoder's $F$ actually is a
contraction is not something proved here or provable by this kind of
argument at all -- it is an empirical property of a specific trained network,
measured per model as $\hat q$ in the main paper's Section~4 and
Section~\ref{sec:s4} below. Symmetrically, Lemma~B (attention dilution) is
unconditional -- it holds for any bounded-spread logits, which is a directly
measurable, not assumed, property of a forward pass -- but its role is only to
license the informal step from ``attention is diluted'' to ``the
per-repetition map is approximately the same map at every boundary''
(Assumption~1 of the main paper), and that step is itself not formalised; see
Section~\ref{sec:s2}.

% ============================================================
\section{What is deliberately not formalized}
\label{sec:s2}

The main paper's Assumption~1 (periodic conditioning: there is a single map
$F$ with $s_{m+1}=F(s_m)$ for every boundary $m$) is presented as ``what
Lemma~B buys us.'' This is true informally and is stated as such in both the
paper and the \texttt{AttentionDilution.lean} header: \emph{``The link is
quantitative but informal (an $O(\delta+1/k)$ perturbation of the conditioning
yields an $O(\delta+1/k)$ perturbation of $F$); the paper states it as a
remark, not as a formalized theorem.''} This section says precisely what that
remark asserts, why it was not formalised, and what a full formal treatment
would need to add. Nothing here is proved in Lean; this section is pen and
paper by design, and we are explicit about that so the boundary between the
machine-checked and the informal parts of the argument is never ambiguous to
a reader.

\subsection{What the remark actually claims}

Lemma~B, as formalised, is a statement about a single decoder step's
attention profile over the $k$ text spans occupied by $k$ repetitions: if
those spans' logits lie within spread $\delta$, the resulting softmax weights
are within $(e^\delta-e^{-\delta})/k$ of uniform (\texttt{softmax\_dist\_le})
and the profile's entropy is at least $\log k - \delta$
(\texttt{entropy\_ge}). Neither of these is yet a statement about $F$, the
map \emph{between repetition boundaries} that Theorem~A's contraction
hypothesis is about. The step from one to the other requires two things
Lemma~B does not supply:

\begin{enumerate}
\item A model of \emph{what the decoder update depends on}. $F$, as used in
  \texttt{CountingCollapse.lean}, is an autonomous map $S\to S$ with no
  explicit dependence on anything besides the current state $s$. In reality
  the per-repetition update is a function of the current state \emph{and} the
  conditioning context the decoder attends to at that step -- write this as
  $F_c : S \to S$ for a context $c$ drawn from some space of possible
  conditioning contexts (e.g.\ the vector of attended text values, not merely
  the attention weights Lemma~B bounds). Autonomy of $F$ is only recovered
  \emph{if} the sequence of contexts $c_0, c_1, \dots$ actually visited at
  successive repetition boundaries is (nearly) constant.
\item A translation of Lemma~B's bound on \emph{attention weights} into a
  bound on the resulting \emph{context vector}. Lemma~B bounds
  $|\mathrm{softmax}(z)_i - \mathrm{softmax}(z)_j|$, a statement purely about
  the weights. The context the decoder actually receives is a
  weight-averaged combination of \emph{value} vectors, and Lemma~B says
  nothing about those values -- it only constrains how uniformly the weights
  spread across positions whose values are assumed, but not shown, to be
  similar because the $k$ repeated spans carry near-identical text.
\end{enumerate}

Granting both informally -- flat attention over near-identical spans yields a
near-constant context vector, to an error the paper writes as
$O(\delta+1/k)$ -- the remark's claim is: a family of maps $\{F_c\}$ that is
uniformly Lipschitz in $c$ (i.e.\ nearby contexts produce nearby updates,
$d(F_c(s), F_{c'}(s)) \le \kappa\, d_C(c,c')$ for some constant $\kappa$ and
every $s$), fed a sequence of contexts $c_0,c_1,\dots$ that are all within
$O(\delta+1/k)$ of some reference context $c^\ast$, behaves \emph{approximately}
like the single autonomous map $F_{c^\ast}$, with an error that itself scales
as $O(\delta+1/k)$ per step. This is the sense in which Assumption~1 is
``bought'' by Lemma~B: not literally, but up to a perturbation that shrinks as
the repeated block grows.

\subsection{Why this was not formalised}

Formalising the remark above is a strictly larger undertaking than
formalising Theorem~A and Lemma~B as stated, for reasons that are worth being
concrete about rather than hand-waving as ``future work'':

\begin{enumerate}
\item \textbf{A new object.} The current development has no type for
  ``conditioning context'' or for a context-indexed family of maps
  $\{F_c\}_{c\in C}$; this would need to be introduced, along with a metric
  $d_C$ on contexts and a joint-Lipschitz hypothesis relating $d_C$ to the
  induced variation in $F_c$. None of this exists in either Lean file today.
\item \textbf{A quantitative re-derivation in the context metric, not the
  weight metric.} \texttt{softmax\_dist\_le} bounds a difference of scalar
  attention weights. Bounding the resulting difference in \emph{context
  vectors} $d_C(c_i, c_j)$ requires an additional Lipschitz-type assumption on
  the value vectors being attended to (that repeated spans carry not just
  diluted attention but near-identical \emph{values}), which is a second,
  currently unstated and unmeasured premise, and then a short but genuine
  argument (a weighted-average version of the triangle inequality) to convert
  a bound on weight differences plus a bound on value differences into a
  bound on the resulting weighted sum.
\item \textbf{A perturbed-contraction (shadowing) lemma.} Even granting (1)
  and (2), Theorem~A's proof chain (Section~\ref{sec:s1}) is stated for a
  single autonomous $F$. Propagating a per-step drift of size $\rho =
  O(\delta+1/k)$ through it requires a genuinely different fixed-point
  argument: something in the spirit of ``a sequence of $q$-contractions
  $F_{c_0}, F_{c_1}, \dots$ whose contexts all lie within $\rho$ of a common
  reference point produces an orbit that shadows the autonomous orbit of
  $F_{c^\ast}$ up to an additive error that itself stabilises geometrically,
  to a limiting value on the order of $\kappa\rho/(1-q)$.'' This is a standard
  \emph{shape} of result in the perturbation theory of dynamical systems (a
  discrete analogue of structural stability / shadowing for contractions),
  but it is not a lemma mathlib currently provides in a directly reusable
  form for this setting, and instantiating or proving it was out of scope for
  this artifact.
\item \textbf{Propagating the correction through the horizon.} Granting (3),
  every inequality in Theorem~A(i)--(iv) would pick up an additive term of
  order $\kappa\rho/(1-q)$ on top of the geometric $Cq^m$ decay, and the
  horizon of Equation~(1) would gain a $k$-dependent correction. Because
  $\rho = O(\delta+1/k)$ shrinks precisely as $k$ (and hence $m$, the quantity
  the horizon bounds) grows, the correction is plausibly a lower-order effect
  in the regime the paper cares about -- but showing this rigorously is a
  two-scale argument (the perturbation shrinking in the same variable the
  main bound is stated over) that requires genuinely new mathematical content,
  not routine bookkeeping, and was not attempted.
\end{enumerate}

\subsection{What a full treatment would need}

Concretely, extending the formal artifact to close this gap would require, in
order: (a) a formal definition of a conditioning-context space $C$ with a
metric $d_C$ and a context-indexed family $F : C \to S \to S$; (b) a
hypothesis package bundling (i) $\kappa$-Lipschitz continuity of $F$ in its
context argument, uniformly in the state argument, and (ii) a quantitative
bound, in $d_C$, on the distance between the contexts realised at successive
repetition boundaries -- the latter to be derived from a strengthened version
of Lemma~B that bounds the attended \emph{value} vectors' contribution, not
only the attention weights; (c) a perturbed-orbit / shadowing lemma for
metric-space contractions under a bounded per-step context drift, proved or
sourced from mathlib; and (d) a re-derivation of Theorem~A(i)--(iv) with the
additive $O(\kappa\rho/(1-q))$ correction carried through explicitly, together
with an argument (formal or at least made quantitatively precise) that this
correction is dominated by the main geometric term in the $k$-range where the
horizon is informative. None of (a)--(d) is present in the current artifact;
the current Lean development proves Theorem~A and Lemma~B exactly as stated
in Section~\ref{sec:s1}, and the bridge between them remains, honestly, a
remark.

% ============================================================
\section{The benchmark in full}
\label{sec:s3}

The benchmark is built by \texttt{data/stimuli/make\_stimuli.py} and written
to \texttt{data/stimuli/stimuli.jsonl}, one JSON object per line. It contains
$180$ items across five families, confirmed by direct count against the
generated file:

\begin{table}[h]
\centering
\caption{Stimulus families.}
\label{tab:families}
\begin{tabular}{lrl}
\toprule
Family & Count & What it tests \\
\midrule
\texttt{word\_rep} & 60 & a single word repeated $k$ times in a fixed carrier \\
\texttt{control\_word} & 54 & the same carrier, $k$ \emph{distinct} filler words (no repetition) \\
\texttt{sentence\_rep} & 32 & a whole short sentence repeated $k$ times \\
\texttt{twisters} & 24 & a tongue-twister sentence (with internal repetition) repeated $k$ times \\
\texttt{numbers} & 10 & English number phrases with intrinsic digit-word repetition \\
\midrule
Total & 180 & \\
\bottomrule
\end{tabular}
\end{table}

Every item carries an explicit \texttt{item\_id}, \texttt{family},
\texttt{template}, \texttt{text}, \texttt{target\_unit}, \texttt{k}, and
\texttt{expected\_count} -- a ground-truth number of occurrences of the target
unit, so a generation can be scored by counting rather than by string match
against a reference transcript. \texttt{control\_word} items additionally
carry \texttt{control\_of} (the \texttt{item\_id} of the repeated item they
are matched to) and \texttt{control\_units} (the $k$ distinct fillers actually
substituted).

\subsection{The $k$-ladder}

\texttt{word\_rep} and \texttt{control\_word} items are built at
$$k \in \{1,2,3,4,6,8,12,16,24,32\}$$
(\texttt{K\_LADDER} in \texttt{make\_stimuli.py}). The ladder is dense at the
low end and log-spaced above; the source comment explains why directly: the
collapse point $k^\ast$ is expected in the $4$--$16$ range, so that is where
resolution is spent, without paying the instrumentation and generation cost
of testing every integer $k$ up to $32$. \texttt{sentence\_rep} uses a
truncated ladder $k\in\{1,2,3,4,6,8,12,16\}$ (\texttt{SENTENCE\_K}), since a
full sentence repeated $32$ times would be prohibitively long to generate and
transcribe. \texttt{twisters} use only $k\in\{1,2,4\}$, since a twister
sentence already contains internal repetition (Table~\ref{tab:twisters}) and
the family exists to probe alliterative/phonological stress rather than
long-range counting. \texttt{numbers} items are not built on a $k$-ladder at
all: each number phrase is a single fixed string with its own intrinsic
repetition count (Table~\ref{tab:numbers}).

\subsection{Word-repetition carrier templates}

Table~\ref{tab:templates} reproduces all six carrier templates
(\texttt{WORD\_TEMPLATES}) verbatim: the fixed prefix and suffix, the target
word repeated $k$ times between them, and the pool of eight distinct fillers
used to build the matched control at the same $k$. An item's text is
\texttt{prefix + " " + (target repeated k times) + " " + suffix}; a control's
text substitutes \texttt{prefix + " " + (k distinct fillers, cycled through
the pool) + " " + suffix}, guaranteeing the target word itself never appears
in its own control.

\begin{table}[h]
\centering
\caption{Word-repetition carrier templates (\texttt{WORD\_TEMPLATES},
\texttt{data/stimuli/make\_stimuli.py}).}
\label{tab:templates}
\small
\begin{tabular}{@{}lp{2.6cm}p{1.3cm}p{3.0cm}p{5.0cm}@{}}
\toprule
ID & Prefix & Target & Suffix & Filler pool (for the control) \\
\midrule
t1 & The dog was & very & big and it ran across the field. & quite, really, truly, fairly, rather, somewhat, extremely, notably \\
t2 & She said & no & to the offer and then left the room. & yes, maybe, sure, fine, okay, well, right, hmm \\
t3 & He walked & far & into the forest before turning back. & deep, fast, long, wide, high, low, near, past \\
t4 & The light was & blue & before the storm arrived. & green, bright, pale, dim, warm, cold, sharp, soft \\
t5 & The runner moved & quick & along the narrow path. & swift, smooth, steady, light, sharp, loose, tight, clean \\
t6 & They were & really & tired after the long journey. & truly, quite, very, rather, deeply, clearly, plainly, surely \\
\bottomrule
\end{tabular}
\end{table}

\paragraph{Why these six target words.} The source comment states the
selection criterion directly: targets are ``chosen for ASR robustness:
common, monosyllabic-or-disyllabic, no homophone confusions, and unlikely to
be swallowed by Whisper's LM prior.'' Concretely, \emph{very}, \emph{no},
\emph{far}, \emph{blue}, \emph{quick}, and \emph{really} are all common,
short, phonetically distinct from every filler in their own pool, and none
are function words or discourse markers that an ASR language-model prior
might silently elide or merge with an adjacent word -- a real risk for a word
repeated up to $32$ times in a row, which is exactly the regime this
benchmark stresses. This choice matters because the behavioural metric
(Section~\ref{sec:s4}) counts occurrences of the target unit in the ASR
transcript; a target word prone to ASR mis-transcription independent of the
TTS model's behaviour would contaminate that count with ASR noise rather than
TTS behaviour.

\subsection{Number phrases}

Table~\ref{tab:numbers} reproduces all ten number-phrase stimuli
(\texttt{NUMBER\_PHRASES}) verbatim, with the target digit-word and its true
occurrence count. Text is capitalised and given a trailing period at
generation time (\texttt{text.capitalize() + "."}). The family's rationale,
from the source comment: ``English number words are the natural repetition
stress test: the same digit-word recurs at several magnitudes and the model
must keep an internal place-value counter that repetition-attractor dynamics
would destroy.'' Unlike \texttt{word\_rep}, the repeated unit here is not
adjacent tokens but a digit-word recurring at different magnitude positions
(hundreds, thousands, millions), so a correct rendering additionally requires
tracking place value, not merely repeating a token a fixed number of times.

\begin{table}[h]
\centering
\caption{Number-phrase stimuli (\texttt{NUMBER\_PHRASES},
\texttt{data/stimuli/make\_stimuli.py}).}
\label{tab:numbers}
\small
\begin{tabular}{@{}lp{7.8cm}ll@{}}
\toprule
ID & Phrase & Target & Count \\
\midrule
n01 & six hundred sixty six thousand six hundred sixty six & six & 6 \\
n02 & seven hundred seventy seven thousand seven hundred seventy seven & seven & 6 \\
n03 & nine hundred ninety nine thousand nine hundred ninety nine & nine & 6 \\
n04 & three hundred thirty three thousand three hundred thirty three & three & 6 \\
n05 & six hundred sixty six million six hundred sixty six thousand six hundred sixty six & six & 9 \\
n06 & eight hundred eighty eight thousand eight hundred eighty eight & eight & 6 \\
n07 & five hundred fifty five & five & 3 \\
n08 & four hundred forty four & four & 3 \\
n09 & two hundred twenty two million two hundred twenty two thousand two hundred twenty two & two & 9 \\
n10 & one hundred eleven thousand one hundred eleven & one & 5 \\
\bottomrule
\end{tabular}
\end{table}

At scoring time (Section~\ref{sec:s4}), a numeral rendered by Whisper as
digits (e.g.\ \texttt{"666,666"}) is algorithmically expanded back into
English number words by \texttt{expand\_number\_token} in
\texttt{src/common/score\_counts.py}, so a model's output can be counted in
the same normalised token space as every other family; a numeral longer than
$12$ digits (i.e.\ a looping model's runaway output, which Whisper can render
as an absurdly long integer) is deliberately \emph{not} run through
place-value expansion and is instead read off digit-by-digit, on the grounds
that place-value names are undefined past \emph{billion} and a digit-wise
reading is the semantically correct interpretation of what a looping decoder
actually spoke.

\subsection{Sentence repetition and tongue twisters}

\texttt{sentence\_rep} repeats one of four short sentences $k\in\{1,2,3,4,6,8,12,16\}$
times (Table~\ref{tab:sentences}); the scored target is a single content word
within the sentence. \texttt{twisters} instead repeat a whole tongue-twister
sentence $k\in\{1,2,4\}$ times (Table~\ref{tab:twisters}); because several
classic twisters already repeat a word or a phoneme cluster \emph{within one
copy} of the sentence, the ground-truth count is
$\text{expected\_count} = \text{per-copy occurrences} \times k$, not $k$
itself. Two twisters (\texttt{w6}, \texttt{w7}) are alliterative rather than
word-repetitive -- their nominal ``target'' is a short substring
(\texttt{"s"}, \texttt{"fr"}) that occurs in many unrelated words, so the
source sets their per-copy occurrence count to $0$ rather than attempting a
substring count that would be semantically meaningless; these two items
therefore always carry \texttt{expected\_count}$=0$ and exercise
phonological/alliterative rendering stress rather than the counting metric.

\begin{table}[h]
\centering
\caption{Sentence-repetition stimuli (\texttt{SENTENCES},
\texttt{data/stimuli/make\_stimuli.py}). $k \in \{1,2,3,4,6,8,12,16\}$ for
every row.}
\label{tab:sentences}
\begin{tabular}{@{}lll@{}}
\toprule
ID & Sentence & Target \\
\midrule
s1 & The bell rang twice. & bell \\
s2 & He counted the stones. & stones \\
s3 & Rain fell on the roof. & rain \\
s4 & The door stayed open. & door \\
\bottomrule
\end{tabular}
\end{table}

\begin{table}[h]
\centering
\caption{Tongue-twister stimuli (\texttt{TWISTERS},
\texttt{data/stimuli/make\_stimuli.py}). ``Per-copy'' is the number of times
the target occurs within a single copy of the sentence; \texttt{expected\_count}
$=$ per-copy $\times\, k$ for $k\in\{1,2,4\}$.}
\label{tab:twisters}
\small
\begin{tabular}{@{}lp{6.7cm}ll@{}}
\toprule
ID & Sentence & Target & Per-copy \\
\midrule
w1 & She sells sea shells by the sea shore. & sea & 2 \\
w2 & Peter Piper picked a peck of pickled peppers. & peck & 1 \\
w3 & How much wood would a woodchuck chuck. & chuck & 1 \\
w4 & Red lorry yellow lorry red lorry yellow lorry. & lorry & 4 \\
w5 & The sixth sick sheikh's sixth sheep is sick. & sixth & 2 \\
w6 & Six slick slim sycamore saplings. & s (alliteration only) & 0 \\
w7 & Fresh French fried fly fritters. & fr (alliteration only) & 0 \\
w8 & Truly rural truly rural truly rural. & rural & 3 \\
\bottomrule
\end{tabular}
\end{table}

\subsection{The matched-control construction}

Every \texttt{word\_rep} item at $k\ge 2$ has a corresponding
\texttt{control\_word} item ($k=1$ is defined to be its own control, since
there is no repetition to remove): the same carrier prefix and suffix, the
same word count, with the $k$ copies of the target replaced by $k$
\emph{distinct} words drawn from that template's filler pool (cycled and
skipped past any accidental collision with the target itself, so the target
word is guaranteed never to appear in its own control). This is the
construction underlying the paper's central dissociation test
(Section~4.2 of the main paper, Figure~1(b)): the control holds carrier
syntax and word count fixed while removing exactly one variable, periodicity,
so that a difference in outcome between the repeated item and its control
cannot be explained by sequence length, generation duration, or the amount of
text the model has to render.

\subsection{The \texttt{boundary\_units} field}

\texttt{word\_rep} and \texttt{sentence\_rep} items carry
\texttt{boundary\_units}: the ordered list of the $k$ words (all copies of
the target, for a repeated item) whose text positions delimit the $k$
repetition boundaries. \texttt{control\_word} items carry the analogous list
of the $k$ distinct fillers actually substituted (duplicated under the key
\texttt{control\_units} for the separate purpose of scoring: was each of the
$k$ fillers rendered at all). This field exists so that the boundary-location
procedure of \texttt{src/common/boundaries.py}
(\texttt{unit\_columns}, Section~\ref{sec:s5}) is handed an explicit,
model-agnostic list of search targets rather than having to special-case
repeated versus control text -- both item types get their boundaries located
by the identical left-to-right, non-rewinding string-search procedure, over
the identical field, which is what makes the fitted decay rates of repeated
and control items directly comparable and not confounded by using two
different localisation methods for the two arms of the comparison.
\texttt{numbers} and \texttt{twisters} items carry no \texttt{boundary\_units}
field at all: numbers have no single repeating word-level unit (the
repetition is a digit-word recurring at different magnitude positions), and
twisters repeat a whole sentence rather than a word, so neither family has a
boundary structure at the granularity this field is designed for.

\subsection{The instrumented subset}

An item is marked \texttt{instrumented=true} (and therefore receives the
teacher-forced hidden-state/attention capture pass described in
Section~\ref{sec:s6}) exactly when its family is one of
\{\texttt{word\_rep}, \texttt{sentence\_rep}, \texttt{control\_word}\} and
$k\ge 2$. The source comment gives the rationale directly: ``Twisters have no
single repeating unit at $k=1$, numbers are short; both are behavioural-only.''
$k\ge2$ is required because at $k=1$ there is by definition no repetition
boundary to instrument. This yields $54\ (\texttt{word\_rep}, k\ge2) + 54\
(\texttt{control\_word}, \text{all } k\ge2\text{ by construction}) + 28\
(\texttt{sentence\_rep}, k\ge2) = 136$ instrumented items, matching the count
written by \texttt{make\_stimuli.py} at generation time and the
\texttt{NInstrumented} macro emitted by \texttt{analysis/make\_numbers.py}.

% ============================================================
\section{Measurement protocol in full}
\label{sec:s4}

\subsection{The conjunctive correctness criterion}

The central behavioural question -- given text asking for $k$ repetitions, how
many did the model actually produce -- is answered by reconciling two
independent, differently-biased estimators, implemented in
\texttt{src/common/score\_counts.py}.

\textbf{Count A (transcript-based).} The generated audio is transcribed with
Whisper large-v3 (\texttt{src/common/asr\_transcribe.py}), the transcript is
normalised (lower-cased, punctuation stripped except commas needed for
numeral parsing, numerals expanded back to number words via
\texttt{expand\_number\_token}), and \texttt{count\_occurrences} counts exact
whole-token matches of the target unit (or, for multi-word targets, greedy
non-overlapping subsequence matches). Count A is accurate when Whisper
faithfully renders every repetition, which holds reliably at low $k$ and
\emph{unreliably at high $k$} -- Whisper is itself an autoregressive model and
is prone to de-duplicating or collapsing long runs of near-identical speech,
exactly the failure mode under study, so it cannot be assumed neutral at the
$k$ values that matter most.

\textbf{Count B (duration-based).} Per $(\text{model}, \text{family},
\text{template})$ cell, a linear fit $\text{duration} = a + b\cdot k$ is
estimated on the \emph{trusted low-$k$ regime}, defined \emph{before} looking
at any high-$k$ behaviour so the calibration cannot be contaminated by the
effect being measured: items with $k\le 3$, Count A exactly equal to the
requested $k$, and audio duration $>0.2\,$s
(\texttt{fit\_duration\_models}). The fit falls back from
$(\text{model},\text{family},\text{template})$ to $(\text{model},\text{family})$
to $(\text{model})$ when a cell has too few trusted points, and is discarded
if the fitted slope $b$ is not meaningfully positive. Inverting the fit at a
given observed duration gives Count B, a count that is \emph{immune} to ASR
de-duplication (it never looks at the transcript) but is correspondingly
\emph{blind} to a model that babbles for approximately the right length
without actually saying the right thing.

Because the two estimators fail in different, largely uncorrelated
directions, \textbf{correctness is decided conjunctively}, not by
picking one estimator or averaging them: an item is labelled \texttt{correct}
only if Count A exactly equals the expected count \emph{and} the observed
duration falls within a tolerance band of the duration the calibrated fit
predicts for that expected count
($\text{duration\_ratio} \in [0.70, 1.45]$, constants \texttt{DUR\_LO},
\texttt{DUR\_HI}). The band is wide enough to absorb ordinary prosodic
speech-rate variation (a model reading the same words slightly faster or
slower is not a counting failure) while still catching the two directions of
genuine failure: a transcript that Whisper collapsed but whose audio ran
much longer than it should have (loop), and a transcript that is short and
whose audio is also short (truncation). Neither estimator is ever used alone
to \emph{assert} a count in a disagreement case -- the reconciliation only
ever \emph{certifies} correctness when both streams of evidence agree; a
disagreement produces an \texttt{agree=False} audit flag
($|\text{Count A} - \text{Count B}| > \max(1, 0.15\cdot\text{expected\_count})$)
rather than a synthesised number, so that a wrong count can never silently
enter the behavioural CSV disguised as ground truth.

\subsection{The outcome taxonomy}

\texttt{classify} in \texttt{score\_counts.py} assigns exactly one of eight
outcome labels, checked in the following fixed order (audio-level degeneracy
first, since noise or silence of coincidentally the right duration would
otherwise slip through the duration test):

\begin{enumerate}
\item \texttt{empty} -- duration $<0.25\,$s or RMS $<10^{-3}$: essentially no
  audio was produced.
\item \texttt{degenerate} -- spectral flatness $>0.35$ (noise/babble rather
  than speech) or an empty ASR transcript.
\item \texttt{correct} -- Count A equals the expected count \emph{and} the
  duration ratio is in-band (the conjunctive criterion above).
\item \texttt{loop} -- generation hit the model's token cap, \emph{or} the
  duration ratio exceeds $1.6$: the model kept going well past where it
  should have stopped.
\item \texttt{overcount} -- Count A exceeds the expected count (and the item
  was not already classified as a loop).
\item \texttt{truncation} -- the duration ratio is below \texttt{DUR\_LO}:
  the model stopped early.
\item \texttt{undercount} -- Count A is below the expected count (and the
  item was not already classified as a truncation).
\item \texttt{miscount} -- the residual fallback: none of the above applied,
  e.g.\ the transcript count matches but the duration ratio is out of band
  without meeting either the loop or the truncation threshold (``right count,
  implausible pacing'').
\end{enumerate}

The main paper's Section~4.3 describes this taxonomy at the six-way level a
figure caption can carry --
$\{\texttt{correct}, \texttt{undercount}, \texttt{truncation},
\texttt{overcount}, \texttt{loop}, \texttt{degenerate}\}$ -- which folds
\texttt{empty} into \texttt{degenerate} conceptually (both are audio-level
failures rather than counting failures) and omits the low-frequency
\texttt{miscount} residual bucket; the eight-way taxonomy above is the one
actually implemented and emitted into \texttt{data/results/behavioural.csv}'s
\texttt{outcome} column.

\subsection{Audio degeneracy detectors}

Computed by \texttt{audio\_flags} in \texttt{src/common/asr\_transcribe.py},
independent of the ASR transcript, because a TTS failure does not always
manifest as wrong \emph{words} -- it can produce noise, a stuck vowel, or
silence, none of which a transcript-only metric would catch:

\begin{itemize}
\item \textbf{RMS} -- root-mean-square amplitude of the waveform; near-zero
  flags near-silence.
\item \textbf{Voiced fraction / trailing silence} -- a smoothed energy
  envelope (frame length $\mathrm{sr}/50$) is thresholded at $2\%$ of its
  peak to flag voiced frames; \texttt{trailing\_silence\_s} is the time from
  the last voiced frame to the end of the clip, distinguishing an abrupt
  truncation (long trailing silence relative to content) from a clean
  stop.
\item \textbf{Spectral flatness} -- \texttt{librosa}'s spectral flatness
  measure, near $1$ for noise-like spectra and near $0$ for tonal/harmonic
  (speech-like) spectra; used as the \texttt{degenerate} threshold above.
\item \textbf{Loop autocorrelation} -- the log-mel envelope (32 mel bands,
  hop length 256, mean over frequency) is autocorrelated with itself; a
  strong peak (\texttt{loop\_ac\_peak}) at a lag beyond $0.4\,$s
  (\texttt{min\_lag}) indicates the model is re-rendering materially the same
  acoustic content periodically -- an exact-loop signature distinct from,
  and complementary to, the duration-ratio-based loop detection in
  \texttt{classify}.
\end{itemize}

Whisper transcription itself is run with
\texttt{condition\_on\_prev\_tokens=False}. This is deliberate, not a default
left in place: Whisper's own decoder is autoregressive and prone to the same
class of repetition lock-in under study, and allowing it to condition on its
own previously decoded text would risk the ASR judge's loops contaminating
the measurement of the TTS model's loops. Word-level timestamps
(\texttt{return\_timestamps="word"}) are requested for a specific downstream
purpose beyond transcription: mapping an audio-domain repetition boundary
back to a decoder step index via $\text{step} = t_{\text{word}} \times
\text{token\_rate\_hz}$, where \texttt{token\_rate\_hz} is each model's
nominal acoustic-token rate from the panel registry
(\texttt{src/common/registry.py}: $50.0\,$Hz for the Llasa family,
$21.53\,$Hz for XTTS-v2, $12.5\,$Hz for Qwen3-TTS). The registry's own
docstring notes that nominal rates are cross-checked against the measured
$n_{\text{tokens}}/\text{duration}$ at smoke-test time, and it is the
measured value that downstream analysis actually uses.

\subsection{Effective rank, spectral slope, and Kirchhoff index}

Computed by \texttt{spectral\_proxies} in \texttt{analysis/state\_dynamics.py}
on a trajectory $h \in \mathbb{R}^{T\times d}$ of hidden states at one probe
layer (mean-centred over time before the decomposition; trajectories longer
than \texttt{MAX\_SVD\_STEPS}$=512$ steps are uniformly subsampled first,
since these are properties of the \emph{shape} of the spectrum, which
uniform subsampling preserves, and the alternative is an $O(Td^2)$ SVD per
layer per item; trajectories shorter than $24$ steps are skipped). Let
$\sigma_1\ge\sigma_2\ge\cdots$ be the singular values of the centred
trajectory matrix and $p_i = \sigma_i/\sum_j\sigma_j$.

\begin{itemize}
\item \textbf{Effective rank} $\mathcal{N}_{\mathrm{eff}} =
  \exp\!\big(-\sum_i p_i\log p_i\big)$: the exponentiated Shannon entropy of
  the normalised singular-value spectrum, i.e.\ a participation-ratio-style
  count of how many singular directions are actually being used, smoothly
  interpolating between $1$ (all variance in one direction) and the nominal
  matrix rank (perfectly flat spectrum). This is the statistic behind the
  capacity-gain measurement of Section~4.3 of the main paper.
\item \textbf{Spectral slope} $\alpha$: the negative slope of $\log\sigma_i$
  regressed against $\log(i+1)$ over the mid-tail index window
  $i \in [2, \max(6, 0.8\cdot\text{rank})]$ -- excluding the first couple of
  dominant components and the extreme tail -- i.e.\ the exponent of the
  spectrum's power-law decay.
\item \textbf{Kirchhoff index} $\log_{10} Kf = \log_{10}\!\big(\sum_i
  1/\lambda_i\big)$ with $\lambda_i=\sigma_i^2$, floored at
  $\lambda_i > 10^{-10}\lambda_{\max}$ for numerical stability: the standard
  graph-theoretic Kirchhoff index (sum of inverse Laplacian eigenvalues, a
  resistance-distance-based measure of global connectivity) applied to the
  covariance spectrum of the trajectory, carried over unchanged from the
  companion Whisper study~\cite{viakhirev2026dispersion} so its Table~2 and
  the present study's Table~2 are directly comparable.
\end{itemize}

\subsection{The template bootstrap}

Confidence intervals throughout \texttt{analysis/capacity.py} are resampled
over \emph{stimulus templates}, not over items, and the rationale is stated
directly in the source: items sharing a template (e.g.\ every $k$-value built
from carrier \texttt{t1}, ``The dog was \dots\ big and it ran across the
field.'') are not independent draws, since they share carrier syntax and
length structure; bootstrapping over raw items would understate the true
variance. Two statistics are computed this way.

\textbf{Capacity gain} (\texttt{gain}): for a given family (repeated or
control) and model, $\mathrm{d}\mathcal{N}_{\mathrm{eff}}/\mathrm{d}\log k$ is
the ordinary-least-squares slope of $\mathcal{N}_{\mathrm{eff}}$ against
$\log k$ at deep probe layers (\texttt{layer\_frac}$>0.6$), fit on items with
$k\ge2$, requiring at least four distinct $k$ values and eight points to
attempt a fit at all. The $95\%$ interval is built by resampling the set of
templates with replacement $2000$ times, refitting the slope on the union of
each resample's items every time. A ``saturation'' value is also reported: the
median $\mathcal{N}_{\mathrm{eff}}$ over items in the top half of the
observed $k$ range, as a plateau estimate.

\textbf{Paired gap} (\texttt{paired\_gap}): for every (template, $k$) cell
with both a repeated-item median and a control-item median available, the
paired difference (control minus repeated) is computed; because the stimuli
were built as matched pairs, this differencing cancels template identity and
$k$ exactly, making it the estimator with the least nuisance variance among
those available. The $95\%$ interval bootstraps the \emph{mean} of the vector
of paired differences ($2000$ resamples with replacement), and
$\texttt{frac\_pos}$ reports the fraction of paired differences that are
positive, i.e.\ the fraction of (template, $k$) cells where the control
gained strictly more distinguishable states than its matched repeated twin.

% ============================================================
\section{The negative result, in detail}
\label{sec:s5}

This section expands the main paper's Section~4.4 (``A note on what we could
not measure''). We consider it a real contribution of the project, not a
discarded attempt, and document it at the level of detail a reviewer wanting
to check the diagnosis would need. The relevant code is
\texttt{src/common/boundaries.py} (explicitly marked, in its own module
docstring, as \emph{``SUPERSEDED as the paper's primary measurement''}) and
\texttt{src/common/dispersion.py} (the estimator that replaced it, marked
\emph{``the primary state measurement''}).

\subsection{What the boundary-difference estimator was}

Theorem~A is stated about the sequence of \emph{repetition-boundary states}
$s_0, s_1, s_2, \dots$ -- the decoder's hidden state exactly at the moment
each successive copy of the repeated phrase has been rendered. The most
direct way to test the contraction premise is therefore to locate those
states in a recorded generation trajectory and difference them. The
estimator built to do this, in order:

\begin{enumerate}
\item \textbf{Locate the $k$ ground-truth text columns.} For a repeated item
  with target unit $u$ (or a control with $k$ distinct fillers),
  \texttt{unit\_columns} walks the model's own tokenized, transformed prompt
  text left to right with a non-rewinding regular-expression search
  (\texttt{\textbackslash b}\texttt{word}\texttt{\textbackslash b}, case-insensitive), using the
  model-specific \texttt{OffsetTokenizer} (\texttt{src/common/offset\_tok.py})
  so that each family's own string transform is respected -- notably XTTS-v2,
  which lower-cases and replaces every space with a literal
  \texttt{[SPACE]} token before tokenizing, so column indices must be (and
  are) computed on the transformed string, not the raw one. This gives $k$
  exact text-span column indices, one per occurrence, by construction in the
  same left-to-right order for a repeated item and its matched control.
\item \textbf{Track an attention centroid.} At the deepest instrumented probe
  layer, \texttt{read\_head} computes, for every generation step $t$, the
  row-normalised text-attention centroid $c_t = \sum_j j\, p(t,j)$ (a
  $5$-tap moving average is applied for smoothing). A centroid rather than an
  arg\-max is used deliberately: the source comment notes that when the $k$
  spans are near-identical, their attention columns are near-identical too,
  so a per-column argmax collapses to the same step for every occurrence,
  whereas a centroid degrades gracefully and still moves even when no single
  column dominates.
\item \textbf{Match columns to steps monotonically.}
  \texttt{boundary\_steps\_from\_attention} rescales the $k$ target columns
  onto the empirical $2$nd--$98$th percentile range the centroid trace
  actually traverses (a centroid, being a weighted mean, is pulled toward the
  middle of the attended mass and so compresses relative to the raw column
  range), then walks forward through the centroid trace once, for each
  rescaled target in turn finding the first step at which the trace first
  reaches or exceeds it. This is monotone by construction -- it can only ever
  move forward through the generation.
\item \textbf{Difference and fit.} The resulting $k$ boundary steps index
  into the recorded hidden-state trajectory, giving
  $d_m = \lVert h(t_{m+1}) - h(t_m)\rVert$
  (\texttt{boundary\_distances}), and \texttt{fit\_geometric} regresses
  $\log d_m = a + m\log q$ by least squares to recover $\hat q_{\text{bnd}}$
  and an $R^2$.
\end{enumerate}

\subsection{Why it was the natural first choice}

This construction is fully self-contained: it uses only the model's own
recorded attention and the tokenizer's exact character offsets, with no
external ASR pass, no separate alignment model, and no assumption that speech
is rendered at a constant tempo. Unlike the estimator that replaced it
(Section~\ref{sec:s5-replacement}), it operationalises Theorem~A's boundary
states $s_m$ \emph{literally}, at the specific generation step where
occurrence $m$ is actually believed to be rendered, rather than approximating
the theorem's claim through a proxy quantity. If the model's attention
tracked position the way a monotone alignment mechanism would, this would be
the right tool, and it is kept in the codebase (rather than deleted) precisely
because it remains the natural estimator for any future model that does have
a genuinely monotone text read-head.

\subsection{The measurement that killed it}

The estimator's soundness rests on one load-bearing assumption: that the
attention centroid $c_t$ moves \emph{monotonically forward} through the text
as generation proceeds, i.e.\ that it is tracking position, not merely
fluctuating around some fixed point of mass. This is directly checkable by
computing, over consecutive generation steps, the fraction of steps at which
the centroid actually advances ($c_{t+1} > c_t$). A genuine left-to-right
sweep would produce a fraction close to $1$; a process with no net
directional drift -- a random walk -- would produce a fraction close to
$1/2$.

The measured value, reported in the project's running notes
(\texttt{findings.md}) and in the docstrings of both
\texttt{dispersion.py} and \texttt{boundaries.py}, is that \textbf{the
deep-layer attention centroid advances on approximately $51\%$ of generation
steps} -- statistically indistinguishable from an unbiased random walk, and
nowhere near the sweeping behaviour the estimator needs. (We note for
precision that this figure is recorded as an empirical summary in the
project's notes and source comments rather than reproduced by a single
standalone script in the current snapshot of \texttt{analysis/}; it is the
stated empirical justification, cited three times independently across the
codebase, for abandoning the estimator, and it is corroborated by the
concrete symptom described next, which \emph{is} directly reproducible from
committed data.)

That concrete symptom is visible in \texttt{data/results/pooled\_q.csv}
(produced by \texttt{analysis/pooled\_q.py}, which pools boundary-distance
curves across items within a (model, family) cell and regresses
$\log(d_m/d_0)$ on $m$ through the origin, bootstrapping the item set for a
confidence interval). Under the boundary/attention method, the pooled fit for
Llasa-1B recovers $\hat q$ statistically indistinguishable from $1$ in every
instrumented family -- \texttt{control\_word}: $\hat q \approx 1.00$ ($95\%$
CI $[0.99, 1.02]$); \texttt{sentence\_rep}: $\hat q \approx 1.03$ ($95\%$ CI
$[1.00, 1.12]$); \texttt{word\_rep}: $\hat q \approx 0.98$ ($95\%$ CI
$[0.94, 1.02]$) -- i.e.\ \emph{no measurable contraction at all}, for a
model whose boundary-free dispersion estimate (Section~\ref{sec:s5-replacement})
recovers $\hat q$ clearly below $1$ for the same generations. This is a
snapshot from the instrumented panel available at the time of writing and
individual decimal values will move as more of the panel finishes running;
the qualitative pattern -- boundary-method $\hat q \approx 1$ versus
dispersion-method $\hat q < 1$, on the identical underlying trajectories -- is
the diagnosis, not the specific digits.

\subsection{Why the failure is self-defeating by construction}

The failure is not a matter of insufficient smoothing or a fixable
implementation bug; it is a direct structural corollary of Lemma~B
(Section~\ref{sec:s1-dilution}). \texttt{softmax\_dist\_le} proves that the
per-occurrence attention-weight differences shrink as $O(1/k)$, and
\texttt{entropy\_ge} proves the profile's entropy floor grows as $\log k -
\delta$: as the repeated block grows, the attention distribution the
estimator's centroid is computed from becomes provably closer to uniform. A
weighted-mean centroid computed over an increasingly uniform distribution
carries, by the same argument, increasingly little information about which
position is truly being rendered -- in the limit its expected value converges
to the midpoint of the attended span regardless of which occurrence is
actually current. The estimator therefore degrades precisely as $k$ grows,
which is exactly the regime in which a clean boundary estimate would matter
most for testing the theorem. This is why the failure could not have been
patched by using more attention heads, denoising the centroid trace, or any
comparable engineering fix: the very effect the estimator was built to
exploit (identifiable, position-informative attention) is the effect Lemma~B
proves vanishes as the block it is trying to measure grows.

\subsection{The replacement: boundary-free dispersion}
\label{sec:s5-replacement}

Theorem~A(ii) (\texttt{exists\_indistinguishable}, Section~\ref{sec:s1})
licenses a different, boundary-free operationalisation: it says the states
visited late in a repeated generation become mutually indistinguishable as a
\emph{set}, which is a statement about the \emph{spread} of the visited
states and needs no per-occurrence boundary labels at all.
\texttt{src/common/dispersion.py} implements this directly: the trajectory is
tiled into \texttt{N\_WINDOWS}$=8$ equal-length windows (each requiring at
least \texttt{MIN\_WINDOW}$=8$ steps to support an estimate); within each
window, states are $\ell_2$-normalised (direction, not residual-stream
magnitude, is what is compared -- magnitude is close to constant along these
trajectories in any case, roughly $130$ for Llasa-1B, so normalising changes
little but makes the resulting quantity comparable across models with
different activation scales) and the mean pairwise chordal distance
$\sqrt{2-2\cos\theta}$ among normalised states in the window is computed;
$\log\sigma(\text{window})$ is then fit linearly against window index, and
the fitted per-window rate is converted to a per-repetition $\hat q$ by the
ratio of window length to average repetition length ($n_{\text{windows}}/k$).
This sidesteps the localisation problem entirely, at the cost of measuring a
slightly different (but, per Theorem~A(ii), equally licensed) quantity: the
spread of the states visited within a window of the trajectory, rather than
the literal distance between two specific, individually-identified boundary
states. \texttt{analysis/state\_dynamics.py} computes both estimators side by
side for every instrumented item (columns \texttt{q\_hat} for dispersion and
\texttt{q\_bnd} for the boundary method in
\texttt{data/results/state.csv}) and records which localisation method,
\texttt{attention} or the uniform-tempo fallback \texttt{uniform}, produced
each row's boundaries, precisely so this substitution is auditable rather
than silent.

% ============================================================
\section{Per-model implementation notes}
\label{sec:s6}

Table~\ref{tab:panel} reproduces the model panel
(\texttt{src/common/registry.py}) in full; the notes below detail how
instrumentation was obtained for each family and, for Qwen3-TTS, precisely
what could not be.

\begin{table}[h]
\centering
\caption{Model panel (\texttt{PANEL}, \texttt{src/common/registry.py}).}
\label{tab:panel}
\footnotesize
\begin{tabular}{@{}p{1.5cm}p{3.4cm}l p{0.8cm} p{1.3cm} p{1.5cm} p{4.4cm}@{}}
\toprule
Key & HF id & Fam. & Par. & Token rate & Probes & Notes \\
\midrule
llasa1b & HKUSTAudio/\allowbreak Llasa-1B & llasa & 1B & 50.0 Hz & all & Llama-3.2-1B / X-codec2 \\
llasa3b & HKUSTAudio/\allowbreak Llasa-3B & llasa & 3B & 50.0 Hz & every 2 & Llama-3.2-3B, same tokenizer/codec \\
llasa8b & HKUSTAudio/\allowbreak Llasa-8B & llasa & 8B & 50.0 Hz & every 4 & Llama-3.1-8B, top of scale ladder \\
xtts2 & coqui/\allowbreak XTTS-v2 & xtts & 0.4B & 21.53 Hz & every 3 & GPT2-style AR + HiFiGAN \\
qwen06b & Qwen/\allowbreak Qwen3-TTS-\allowbreak 12Hz-\allowbreak 0.6B-Base & qwen & 0.6B & 12.5 Hz & every 3 & dual-track multi-codebook LM \\
qwen17b & Qwen/\allowbreak Qwen3-TTS-\allowbreak 12Hz-\allowbreak 1.7B-Base & qwen & 1.7B & 12.5 Hz & every 3 & same family as 0.6B \\
\bottomrule
\end{tabular}
\end{table}

\texttt{probe\_layer\_indices} always force-includes the final layer of
whichever density setting is used, since that is the layer that actually
feeds the model's own readout (stop-token logit or codec head); the ``every
$n$'' densities on the larger checkpoints exist to bound per-item activation
storage as depth grows.

\subsection{Llasa (1B / 3B / 8B)}

\texttt{src/models/llasa\_gen.py} runs every Llasa checkpoint through the
same two-pass design. \textbf{Pass 1 (generate)} samples speech tokens
autoregressively via \texttt{model.generate} with the attention backend set
to SDPA (\texttt{do\_sample=True}, temperature $0.8$, top-$p$ $1.0$,
\texttt{max\_new\_tokens}$=2048$ by default); the prompt is built from the
model's own chat template with sentinel tokens
(\texttt{<|TEXT\_UNDERSTANDING\_START/END|>},
\texttt{<|SPEECH\_GENERATION\_START/END|>}) and
\texttt{continue\_final\_message=True} to prime the assistant turn on the
speech-start token; generated speech-token ids are extracted from the
\texttt{<|s\_NNN|>} vocabulary pattern and saved as integer arrays, with
waveform synthesis deferred to a separate offline pass
(\texttt{src/models/xcodec2\_decode.py}) run in the isolated
\texttt{xcodec2} environment (Section~\ref{sec:s7}). A \texttt{hit\_cap} flag
records whether generation ran to the token ceiling; per the project's
running notes, $21.5\%$ of Llasa-1B generations do so, treated in this study
as a signal (runaway generation is common enough to be informative) rather
than noise to be tuned away by simply raising the cap.

\textbf{Pass 2 (instrument)} runs only when the seed equals
\texttt{instrument\_seed} (default $0$), the stimulus item is marked
\texttt{instrumented}, and at least one speech token was produced. The
\emph{entire realised sequence} (prompt followed by the tokens actually
sampled in Pass 1, EOS stripped) is fed through a single teacher-forced
forward call with the attention backend switched to \texttt{eager} -- the
only backend that materialises per-layer attention-weight tensors -- and
\texttt{output\_hidden\_states=True}. Because the model is causal, position
$t$ of this forward pass carries exactly the hidden state that \emph{produced}
generated token $t$ during the actual sampling pass, so this single clean
forward recovers the whole instrumented trajectory without altering the
realised token sequence in any way; it is a pure read-out of Pass 1's output,
not a re-sampling. \texttt{TextAttentionRecorder} attaches forward hooks to a
chosen subset of \texttt{self\_attn} layers (up to four of
$\{0, \lfloor L/3\rfloor, \lfloor 2L/3\rfloor, L-1\}$ for $L$ total layers);
each hook slices the returned \texttt{(attn\_output, attn\_weights)} tuple's
weights down to \texttt{[B, H, Q, text\_lo:text\_hi]}, head-averages, and
moves to CPU \texttt{float16} \emph{inside the hook}, specifically to avoid
ever materialising the full $O(T^2\cdot\text{heads}\cdot\text{layers})$
attention tensor in memory. Entropy and top-1 probability of the next-token
distribution are recorded from the same forward pass's logits at no extra
cost. \texttt{llasa\_gen.py}'s own comment states the reason SDPA is used for
sampling and eager only for the instrumented pass directly: Llasa under
eager attention is roughly $5\times$ slower to sample.

\subsection{XTTS-v2}

\texttt{src/models/xtts\_gen.py} generates via \texttt{model.gpt.generate}
(the GPT2-style autoregressive mel/codec-index decoder), using the sampling
defaults shipped in the checkpoint's \texttt{config.json} unless overridden
on the command line. Those shipped defaults include
\textbf{\texttt{repetition\_penalty}$=5.0$} -- confirmed directly in the
downloaded checkpoint's \texttt{config.json} -- an aggressive built-in
anti-repetition mechanism active during the very sampling pass this study
measures. Any residual counting failure XTTS-v2 exhibits is therefore
occurring \emph{despite} a penalty actively working against it, which is why
the main paper's Discussion (Limitations) frames XTTS-v2 as a
\emph{conservative}, not a neutral, member of the panel: if anything, the
theory's prediction is being tested on this model under a handicap. The AR
decoder's built-in generation ceiling
(\texttt{model.gpt.max\_gen\_mel\_tokens}) is derived from the shipped
config's \texttt{model\_args.gpt\_max\_audio\_tokens}$=605$ minus a small
fixed number of reserved special-token positions, giving an effective cap of
approximately $602$ mel frames per item; the script's own
\texttt{--max-new-tokens} override can only shrink this further via
\texttt{min()}, never raise it.

Instrumentation avoids a second forward pass through the high-level API by
hand-reconstructing XTTS's own internal embedding recipe: conditioning
latents, text embeddings (token $+$ positional), and mel-token embeddings
(token $+$ positional) are concatenated into a single \texttt{inputs\_embeds}
tensor and fed through the \emph{raw} underlying \texttt{transformers.GPT2Model}
(\texttt{model.gpt.gpt}, whose own \texttt{wte}/\texttt{wpe} embedding tables
are deleted by XTTS since embeddings are computed externally). This bypasses
\texttt{Xtts.inference()}/\texttt{synthesize()} specifically because those
high-level entry points do not expose \texttt{output\_hidden\_states=True}
and additionally run a sentence-splitting/chunking stage that could silently
fragment a long repeated-word stimulus into several independent
\texttt{generate()} calls, defeating the ``one contiguous generation''
premise the state-dynamics measurement depends on. Attention is forced to
\texttt{eager} once at load time
(\texttt{model.gpt.gpt.config.\_attn\_implementation = "eager"}), because
SDPA (the library default) silently returns \texttt{None} attention weights
even when \texttt{output\_attentions=True} is requested -- a failure mode the
script's header comment documents explicitly, verified against the installed
\texttt{coqui-tts}~0.27.5 and \texttt{transformers}~4.57.1 sources by file
and function name. The same memory-avoidance trick as Llasa is used
(\texttt{output\_attentions=False} at the top-level call, so
\texttt{GPT2Model} never collects the full attention stack across all
layers, while per-layer forward hooks on the probed blocks' \texttt{.attn}
submodules see the real per-layer weights regardless, since
\texttt{transformers}' eager attention path computes them unconditionally
internally). Vocoding runs in the same process and script as generation
(unlike Llasa, which vocodes offline in a separate environment), since
XTTS ships its own HiFiGAN decoder with no conflicting pinned dependency.

\subsection{Qwen3-TTS (0.6B / 1.7B)}

\texttt{src/models/qwen\_gen.py} instruments a model that is architecturally
very different from the other two families. Qwen3-TTS is a dual-track model:
an autoregressive ``talker'' transformer
(\texttt{Qwen3TTSTalkerForConditionalGeneration}) emits one $12.5$\,Hz
acoustic frame per decode step, where each frame packs $16$ RVQ codebook ids
-- the first sampled directly by the talker's own \texttt{codec\_head} inside
the standard HF generation loop, the remaining $15$ filled in afterward by a
separate, smaller ``code predictor'' sub-model conditioned on that step's
talker hidden state. Only the first-codebook / talker level is instrumented,
which is also the level at which Theorem~A's boundary-state argument is
naturally stated.

\textbf{No second forward pass is needed.} The mid-level library API,
\texttt{Qwen3TTSForConditionalGeneration.generate()}, already unconditionally
calls the talker's own \texttt{.generate(...,}
\texttt{output\_hidden\_states=True,} \texttt{return\_dict\_in\_generate=True)}
internally on every call and then discards every hidden-state layer but the
last. \texttt{TalkerGenerateCapture} monkeypatches the bound
\texttt{talker.generate} method to stash the full raw
\texttt{GenerateDecoderOnlyOutput} (and the \texttt{trailing\_text\_hidden}
keyword argument) before forwarding the call through completely unchanged --
a read-only interception, so sampling and behaviour are byte-for-byte
identical to what the public \texttt{generate\_voice\_clone()} API would have
produced with or without instrumentation. Per-frame, per-probe-layer hidden
states are then pulled out of the captured object after the fact
(\texttt{extract\_trajectory}), with frame $j$ aligned to the decode call
whose forward pass predicted it, matching the library's own internal
step/frame convention.

\textbf{Text attention is not obtainable, and this is stated in the code as
an architectural fact, not a shortfall of instrumentation effort.} In the
model's default streaming mode, the embedding fed into the talker at each
decode step is the \emph{elementwise sum} of a codec-token embedding and, for
a short prefix of steps, a text-token embedding
(\texttt{trailing\_text\_hidden}, added inside the talker's own
\texttt{forward}). There is consequently no column of the self-attention
matrix that corresponds to ``the text'' the way there is for Llasa's or
XTTS's concatenated token streams: attention weight over a fused position
cannot be attributed to its text component versus its codec component after
the sum, because the sum is not invertible. A \texttt{non\_streaming\_mode=True}
flag exists in the library that would give text its own contiguous block of
positions and could in principle be probed, but it is not the model's
default or tested code path, and using it risks degrading the audio quality
the project treats as the non-negotiable baseline requirement, so it is
deliberately not used. \texttt{text\_lo}/\texttt{text\_hi} are therefore
written as literal \texttt{null} in this model's metadata records -- an
explicit ``not applicable,'' never a fabricated or uniform placeholder span
standing in for real boundaries -- alongside a \texttt{trailing\_text\_len}
field recording how many initial decode steps did carry a distinguishable
text contribution, offered as a temporal (not spatial) proxy for anyone who
wants one. A direct consequence, made explicit here because it is easy to
miss: because Qwen3-TTS has no obtainable text attention, it is absent from
the boundary-based estimator of Section~\ref{sec:s5} \emph{by construction}
(there is no attention tensor to feed \texttt{unit\_columns}/\texttt{read\_head}),
and its state-space measurements in this study are dispersion-only.

% ============================================================
\section{Reproducibility}
\label{sec:s7}

\subsection{The four conda environments}

\texttt{env/setup\_envs.sh} builds four environments because, per its own
header comment, ``the dependency constraints are mutually unsatisfiable.''
Table~\ref{tab:envs} lists every version pin found in the script and the
concrete failure it exists to prevent, drawn from the script's own comments
and the README's ``Things that will bite you'' section.

\begin{table}[h]
\centering
\caption{Conda environments (\texttt{env/setup\_envs.sh}) and their pins.}
\label{tab:envs}
\small
\begin{tabular}{@{}lp{2.6cm}p{8.6cm}@{}}
\toprule
Env & Python / core & Pin and the failure it prevents \\
\midrule
\texttt{base} & py3.13, torch 2.11+cu130 & \texttt{transformers==4.57.3} (additive install onto
  the existing base torch); runs the Llasa LMs, Whisper ASR, and all analysis
  scripts. \\
\texttt{qwen} & py3.12 & \texttt{qwen-tts} package; torch/torchaudio installed from the
  \texttt{cu130} index \emph{before} \texttt{qwen-tts} -- installing in the
  other order lets pip silently pull a default CUDA-12 wheel, breaking the
  driver/toolkit match. \\
\texttt{coqui} & py3.12 & \texttt{coqui-tts}; same cu130-before-package ordering.
  \texttt{coqui-tts} breaks on \texttt{transformers} 5.x
  (\texttt{isin\_mps\_friendly} was removed from that version), so this
  environment is kept on \texttt{transformers} 4.57.x. \\
\texttt{xcodec2} & py3.11, torch 2.5 (hard-pinned) & \texttt{xcodec2==0.1.5} pins torch 2.5 but
  leaves \texttt{transformers}/\texttt{torchao} unpinned, and both have since
  moved past compatible versions: \texttt{transformers} 5.x $\to$
  \texttt{Could not import module 'PreTrainedModel'}; \texttt{torchao}
  $\ge$0.7 $\to$ \texttt{module 'torch' has no attribute 'int1'}. The working
  combination, pinned explicitly after the \texttt{xcodec2} install, is
  \texttt{transformers==4.46.3}, \texttt{tokenizers<0.21},
  \texttt{torchao==0.6.1}, plus \texttt{soundfile}. \\
\bottomrule
\end{tabular}
\end{table}

The \texttt{qwen} and \texttt{coqui} environments both install torch and
torchaudio from the PyTorch \texttt{cu130} wheel index
(\texttt{--index-url} \texttt{https://download.pytorch.org/\allowbreak whl/cu130})
\emph{before} the TTS package itself, per the script's own top-level
comment: ``install cu130 torch BEFORE the TTS package in each env, else pip
silently pulls a default CUDA-12 wheel and the driver/toolkit combo breaks.''

\subsection{Other pitfalls documented in the README}

\begin{itemize}
\item \textbf{\texttt{src/common/tokenizers.py} must never exist.} Running a
  script from inside \texttt{src/common/} puts that directory first on
  \texttt{sys.path}, so a file with that name would shadow the real
  \texttt{tokenizers} package and break \texttt{transformers} with a
  confusing \texttt{ImportError}; this is why the project's own tokenizer
  interface is named \texttt{offset\_tok.py} instead.
\item \textbf{The HF \texttt{automatic-speech-recognition} pipeline routes
  audio through \texttt{torchcodec}}, which does not load against the
  installed FFmpeg in this environment; \texttt{asr\_transcribe.py} drives
  \texttt{WhisperProcessor} and \texttt{WhisperForConditionalGeneration}
  directly instead, avoiding that code path entirely.
\item \textbf{XTTS replaces spaces with literal \texttt{[SPACE]} tokens and
  lower-cases before tokenizing}, so any text-column index computed for XTTS
  must be computed on the transformed string
  (\texttt{offset\_tok.\_xtts\_prepare}), never on the raw stimulus text.
\item \textbf{Llasa under eager attention is roughly $5\times$ slower to
  sample}, which is why \texttt{llasa\_gen.py} samples under SDPA and
  switches to eager only for the single instrumented teacher-forced pass.
\item \textbf{Whisper de-duplicates repeated speech}, so a transcript-only
  repetition count systematically understates loops; this is the direct
  motivation for the conjunctive correctness criterion of
  Section~\ref{sec:s4}.
\end{itemize}

\subsection{Exact commands}

\begin{lstlisting}[language=bash,mathescape=false]
# 0. Environments (once)
bash env/setup_envs.sh all       # base, qwen, coqui, xcodec2
bash scripts/setup_lean.sh       # elan + mathlib cache (CPU, ~15 min)
bash scripts/download_models.sh  # 8 checkpoints, ~74 GB

# 1. Verify the formal artifact
bash scripts/check_lean.sh

# 2. Generate (one model per GPU; resumable by (item_id, seed))
python src/models/llasa_gen.py --model llasa1b --gpu 0 --seeds 0 1 2
# xtts2 needs the `coqui` env and COQUI_TOS_AGREED=1
# qwen06b / qwen17b need the `qwen` env

# 3. Vocode, transcribe, score
bash scripts/run_pipeline.sh <asr_gpu> llasa1b xtts2 ...

# 4. Analyse and build the paper
python analysis/state_dynamics.py --models <models> --out data/results/state.csv
python analysis/capacity.py                       # the central measurement
python analysis/horizon_fit.py                     # behavioural collapse points
python analysis/figures.py
bash paper/build.sh                                # regenerates numbers, compiles
\end{lstlisting}

\texttt{paper/build.sh} regenerates \texttt{numbers.tex} from the result CSVs
before every compile (\texttt{analysis/make\_numbers.py}), so a number in the
compiled PDF cannot drift from the data it was computed from; every
generation script is resumable, skipping any \texttt{(item\_id, seed)} pair
already present in its metadata file, so re-running after an interruption is
always safe. Because generation and analysis for this panel were still
in progress at the time this document was written, several of the
per-model quantities in the main paper's Table~1 and Table~2 are
recomputed automatically the next time \texttt{paper/build.sh} runs against a
more complete \texttt{data/results/} directory; this document has
deliberately favoured describing method over quoting point-in-time numbers
for exactly that reason.

\subsection{The verification gate}

\texttt{scripts/check\_lean.sh} is reproduced in full in
Section~\ref{sec:s1-audit}. Its three checks -- successful \texttt{lake
build}, a textual \texttt{sorry} scan, and the \texttt{\#print axioms} audit
for \texttt{sorryAx} -- are run as a single command,
\texttt{bash scripts/check\_lean.sh}, and constitute the project's gate on
the formal artifact: no result in the main paper or in this document should
be read as resting on a machine-checked guarantee unless this script has
been run against the exact commit being cited and has printed
\texttt{PASS: formal artifact verified}.

\bibliographystyle{IEEEtran}
\bibliography{../refs}

\end{document}
